% ============================================================
% M3 S2 导学件·波1 补位4a 片件 —— 课时17 2.8① 直线与圆锥曲线的位置关系（一）位置关系判定与弦长（42键＝探究37＋评价38—42）
% 生成：M3 成卷轮 S2 波1 补位4a｜2026-09-14｜编译 cwd＝本目录（中文路径禁 \input 绝对路径）
%   xelatex main-true.tex ×2 → 含详解印本；xelatex main-false.tex ×2 → 纯题版（单源双本）
% 输入链（全只读）：题面库 课时17-2.8①压轴综合一.md（题面逐字装配）＋ -答案侧.md（值逐字回嵌）
%   ＋ 定稿/批D-课时17-2.8①压轴综合一.md（详解回嵌源：§七 补产全文区＋件末「导学 G5 补产」）；
%   toolchain＝qp-m3.sty（钉后版 md5 7c3930362be8a0a2bdf21bbf8ac16573，本地挂载副本，破损即停）。
% 母版：课时01-坐标法（结构骨架照抄）；近范模：课时16（波1 臂5 片）；门谱：工作区/_tmpM3S2波1补位4a_0914/门谱/。
% 键账：件内 ansblock 键集 ≡ 题面库片键集 ≡ manifest 键序（42 键，零缺漏零多余）；
%   印面号 1—42＝探究 1—37＋课堂评价 38—42，键↔号映射见 值台账-课时17.json；
%   让位席卷④§2.8.12.1-25 让位不设键不设块（键集豁免登记随 manifest 备注），拓键实幅 01~21。
% 括线判据（派单承重墙，禁放宽）：估高＞8 行（wlen（全角1·ASCII0.5，剔\命令）/23 栏宽字口径，
%   est＝ceil）→ 括线模（{\ansblockgrayfalse 包装}），否则灰底；逐键读数入值台账。
% 值快照：42 键值均与答案侧「值：」行逐字全等（零换算；练11 值内花括号以 \{ \} 印面转写、
%   逐字零漂口径见门谱；− 走 preamble 路由、∪∞∈ 已探针实证零缺字）。
% 印面符号路由（探针实证 0914，注册项）：✓→\ding{51}、✗→\ding{55}（pifont）；⇒→$\Rightarrow$、
%   ⟺→$\iff$、∥、∈、≤、≥、∪、∞、−(U+2212 路由)、–(U+2013 路由) 等详见台账渲染注记。
% ============================================================
\documentclass[fontset=none]{ctexart}
\usepackage{qp-m3}
% —— 印面逐字符号路由补钉（探针实证 0914：书宋缺 U+2212/U+2013，▱ 诸字库皆缺→NSC 子块；
%     禁改 sty 本体保 md5；✓✗ 探针实证全字体链皆缺→pifont \ding 路由） ——
\usepackage{pifont}%
\xeCJKDeclareCharClass{Default}{"2212}%
\xeCJKDeclareCharClass{Default}{"2013}%
\setCJKfamilyfont{nscblk}[Path=C:/提示词/工作区/字替对照-0909/variantF/fonts/,BoldFont=NSC-w600.ttf]{NSC-w500.ttf}%
\xeCJKDeclareSubCJKBlock{nscblk}{"25B1}%
\newcommand{\pxparallelogram}{{\CJKfamily{nscblk}▱}}%
% —— 断行微调（CJK 胶 0.10em＋拉伸 plus 0.30em，课时16 实证校准 0914：窄栏详解长公式链
%     0.08em 吸不住必报 underfull，0.30em 后全数吸收；探针 probe6—8 实证） ——
\xeCJKsetup{CJKglue={\hskip 0.10em plus 0.30em minus 0.05em}}%
% —— 答案标签原子化（拓17-04/16 长值 ansitem 首行 justify 拉伸入 [答案] 标签内胶，
%     印面成「[答　案]」裂距 0914 目验在案；\mbox 壳令标签定宽不参调，样式零改） ——
\renewcommand{\anlabel}[1]{\mbox{#1}}%
% —— 页脚件名（qp-headfoot 复用入口） ——
\renewcommand{\qpzhangming}{第二章\quad 平面解析几何}%
\renewcommand{\qpjianming}{导学件}%

\begin{document}

\zhangtitle{第二章\quad 平面解析几何}
\jietitle{2.8 直线与圆锥曲线的位置关系（一）}
\mubiaoline
\mubiaomu{1}{掌握直线与椭圆、双曲线、抛物线位置关系的判定方法，会联立消元看 \(\Delta\) 判位置、求公共点与弦长，并能检查二次项系数是否退化为零；}
\mubiaomu{2}{会按斜率与位置特例分类讨论含参直线与圆锥曲线公共点的个数，识别平行渐近线、平行对称轴等退化特位，辨析「相切」与「恰有一个公共点」；}
\mubiaomu{3}{熟练运用弦长公式与焦点弦、焦半径结论处理弦长与弦中点问题，体会代数判定与几何直观的结合，发展数学运算与逻辑推理素养.}

\vspace{1.7mm}
\begin{multicols}{2}
\emergencystretch=1em

% ============================================================
% 课前预习（知识点导学填空；零键零答案——预习键位按检查口令④裁定前口径，
% \kongbai{} 空线＋\zhentib 判断题面形，两档等形不泄答）
% ============================================================
\huaxing{课}{前}{预}{习}{知识导学\quad 素养初识}

\zsd{一}{位置关系的判定与公共点}

\tiaomu{1}{判定直线与椭圆、双曲线、抛物线的位置关系，通常用\kongbai{} 的方法：联立消元得 \(ax^{2}+bx+c=0\)（\(a\neq 0\)）时，\(\Delta>0\) 等价于有\kongbai{} 个公共点，\(\Delta=0\) 等价于\kongbai{}，\(\Delta<0\) 等价于\kongbai{}．\zhuzhu{消元后必须检查二次项系数：直线平行于双曲线的渐近线或抛物线的对称轴时方程降为一次，恰有一个公共点但不相切.}}

\tiaomu{2}{「直线与抛物线相切」是「直线与抛物线只有一个公共点」的\kongbai{} 条件；「直线与椭圆只有一个公共点」是「直线与椭圆相切」的\kongbai{} 条件．\zhuzhu{椭圆是封闭曲线，消元不降次，恰一公共点即相切；抛物线、双曲线有降次特位，须单独辨析.}}

\zsd{二}{弦长与焦点弦}

\tiaomu{1}{直线 \(y=kx+b\) 与曲线交于 \(A(x_{1},y_{1})\)、\(B(x_{2},y_{2})\) 两点，则弦长 \(|AB|=\)\kongbai{}；抛物线 \(y^{2}=2px\)（\(p>0\)）过焦点的弦长 \(|AB|=\)\kongbai{}．\zhuzhu{弦长公式对三种曲线通用；焦点弦先想焦半径 \(|PF|=x_{P}+p/2\)（即到准线的距离），再化为坐标和.}}

\zhenhead{判断正误(正确的打\gou{},错误的打\cha{})}

\zhentib{(1)}{直线与抛物线只有一个公共点，则直线与该抛物线相切．}

\zhentib{(2)}{直线与椭圆只有一个公共点，则直线与该椭圆相切．}

\liubai[9mm]{此处书写}

% ============================================================
% 课中探究（考点探究〔例+变式〕＝练/拓槽 37 键；题后紧跟答案＝ansblock 内嵌，
% 值照答案侧逐字，详解承批D回嵌＋印面清洗：〔…〕审计尾注剔除、✓→\ding{51}、∎消、⟺⇒转数学宏；
% 17-16 之定位辨析句按定稿引用须知5 必带）
% ============================================================
\huaxing{课}{中}{探}{究}{考点探究\quad 素养小结}

% ---------- 探究点一 位置判定与相切辨析 ----------
\tjdnr{一}{位置判定与相切辨析}{例\textbf{1}}{简单(知识点一)}{判断直线 \(y=-2x+4\) 与抛物线 \(y^{2}=4x\) 是否有公共点．如有，求出公共点的坐标，如公共点有两个，求出以这两个公共点为端点的线段长．}
\xiexwei{12mm}

{\ansblockgrayfalse
\begin{ansblock}[2章-练-课时17-01]
% ans:2章-练-课时17-01
\ansitem{1}{有公共点；两公共点 (1,2)、(4,−4)；以之为端点的线段长 3√5}
\ansnote{详解}{以 \(y\) 为主元消元（直线已给 \(y\)、抛物线给 \(x=y^{2}/4\)）：\(y=-2\times(y^{2}/4)+4=-y^{2}/2+4\)，两边乘 2 整理得 \(y^{2}+2y-8=0\)（二次项系数 \(1\neq 0\)，故确为二次、不会退化）．\(\Delta=2^{2}+4\times 8=36>0\)，故有两个公共点．\(y=(-2\pm 6)/2\)，即 \(y=2\) 或 \(y=-4\)；代回 \(x=y^{2}/4\) 得 \(x=1\) 或 \(x=4\)．\par\noindent 故公共点为 \((1,2)\) 与 \((4,-4)\)，所求线段长 \(=\sqrt{(4-1)^{2}+(-4-2)^{2}}=\sqrt{9+36}=\sqrt{45}=3\sqrt{5}\)．（回验：\(2=-2\times 1+4\) \ding{51}；\(-4=-2\times 4+4\) \ding{51}；\(1=1\)、\(16=4\times 4\) \ding{51}．）}
\end{ansblock}}

\liB{变式\textbf{1}}{简单(知识点一)}{}{直线 \(y=x+1\) 与椭圆 \(x^{2}+y^{2}/2=1\) 的位置关系是\nobreak(\kongwei)}

\bindopt {\fontsize{10.5pt}{\qpTJgLeadLiti}\selectfont \makebox[\dimexpr0.5\linewidth-1em\relax][l]{A．相离}\makebox[\dimexpr0.5\linewidth-1em\relax][l]{B．相切}\par}

\bindopt {\fontsize{10.5pt}{\qpTJgLeadLiti}\selectfont \makebox[\dimexpr0.5\linewidth-1em\relax][l]{C．相交}\makebox[\dimexpr0.5\linewidth-1em\relax][l]{D．无法确定}\par}

\begin{ansblock}[2章-练-课时17-04]
% ans:2章-练-课时17-04
\ansitem{2}{C}
\ansnote{详解}{联立 \(y=x+1\) 与 \(x^{2}+y^{2}/2=1\)，消去 \(y\)：\(x^{2}+(x+1)^{2}/2=1\)，乘 2 得 \(2x^{2}+x^{2}+2x+1=2\)，即 \(3x^{2}+2x-1=0\)．\(\Delta=2^{2}-4\times 3\times(-1)=4+12=16>0\)，且二次项系数 \(3\neq 0\)，故有两个公共点，直线与椭圆相交，选 C．（易错点：消元后须检查二次项系数是否为 0；本题 \(x^{2}\) 项由椭圆方程给出、恒在，无降次之虞．）}
\end{ansblock}

\tjdnr{一}{位置判定与相切辨析}{例\textbf{2}}{简单(知识点一)}{举例说明，直线与抛物线只有一个公共点不是它们相切的充分条件．}
\xiexwei{12mm}

{\ansblockgrayfalse
\begin{ansblock}[2章-练-课时17-02]
% ans:2章-练-课时17-02
\ansitem{3}{反例：直线 y=1 与抛物线 y²=4x 只有一个公共点 (1/4,1)，但二者不相切（该点处抛物线的切线是 y=2x+1/2）}
\ansnote{详解}{取抛物线 \(C\)：\(y^{2}=4x\) 与直线 \(y=1\)（平行于 \(C\) 的对称轴 \(x\) 轴）．把 \(y=1\) 代入得 \(1=4x\)，唯一解 \(x=1/4\)，故恰有一个公共点 \((1/4,1)\)．\par\noindent 再验它不相切：设过该点的直线 \(y=k(x-1/4)+1\)，记 \(t=x-1/4\)（则 \(4x=4t+1\)），代入 \(y^{2}=4x\) 得 \((1+kt)^{2}=1+4t\)，展开整理为 \(k^{2}t^{2}+(2k-4)t=0\)．\(k\neq 0\) 时方程有二根 \(t=0\) 与 \(t=(4-2k)/k^{2}\)，相切须二根重合，即 \((4-2k)/k^{2}=0\) \(\iff k=2\)，切线为 \(y=2(x-1/4)+1\)，即 \(y=2x+1/2\)；而 \(k=0\)（即直线 \(y=1\)）时联立式退化为一次方程 \(-4t=0\)，只有一个公共点却不是重根，按「相切＝联立有重根」的判定，它不是切线、直线与抛物线不相切．\par\noindent 故「直线与抛物线只有一个公共点」成立时相切未必成立，即它不是相切的充分条件．（本席事理与变式 2、例 5 之退化特位三处互证．）}
\end{ansblock}}

\liB{变式\textbf{2}}{简单(知识点一)}{}{「直线与抛物线相切」是「直线与抛物线只有一个公共点」的\nobreak(\kongwei)}

\lxopt{A．充分不必要条件}

\lxopt{B．必要不充分条件}

\lxopt{C．充要条件}

\lxopt{D．既不充分也不必要条件}

\begin{ansblock}[2章-练-课时17-05]
% ans:2章-练-课时17-05
\ansitem{4}{A}
\ansnote{详解}{正向：相切即联立所得二次方程有重根（\(\Delta=0\)），重根只对应一个公共点，故「相切」\(\Rightarrow\)「只有一个公共点」，充分性成立．\par\noindent 反向：抛物线中消元可能降次——直线平行于对称轴时（如 \(y^{2}=4x\) 与 \(y=1\)，见例 2），联立式退化为一次方程，恰有一个公共点却不相切，故「只有一个公共点」未必能推「相切」，必要性不成立．\par\noindent 于是「相切」是「只有一个公共点」的充分不必要条件，选 A．}
\end{ansblock}

\tjdnr{一}{位置判定与相切辨析}{例\textbf{3}}{简单(知识点一)}{判断正误．}

\bindp (1) 过椭圆外一点只能作一条直线与椭圆相切．

\bindp (2) 直线 \(y=x\) 与椭圆 \(x^{2}/a^{2}+y^{2}/b^{2}=1\)（\(a>b>0\)）不一定相交．

\bindp (3) 过点 \((3,0)\) 的直线有且仅有一条与椭圆 \(x^{2}/9+y^{2}/16=1\) 相切．

\bindp (4) 直线与椭圆只有一个交点 \(\iff\) 直线与椭圆相切．
\xiexwei{12mm}

{\ansblockgrayfalse
\begin{ansblock}[2章-练-课时17-07]
% ans:2章-练-课时17-07
\ansitem{5}{(1)错误　(2)错误　(3)正确　(4)正确}
\ansnote{详解}{(1) 错．椭圆外一点可作两条切线：如过 \((4,0)\) 对椭圆 \(x^{2}/4+y^{2}=1\)，设 \(y=k(x-4)\) 代入得 \((1+4k^{2})x^{2}-32k^{2}x+64k^{2}-4=0\)，令 \(\Delta=16-192k^{2}=0\) 得 \(k=\pm\sqrt{3}/6\)，恰两条，故「只能一条」错．\par\noindent (2) 错．原点 \((0,0)\) 在椭圆内部（\(0+0<1\)），过内点的直线必与椭圆交于两点；直接联立 \(b^{2}x^{2}+a^{2}y^{2}=a^{2}b^{2}\) 与 \(y=x\) 得 \((a^{2}+b^{2})x^{2}=a^{2}b^{2}\)，即 \(x^{2}=a^{2}b^{2}/(a^{2}+b^{2})>0\)，确有两交点，故「不一定相交」错．\par\noindent (3) 对．\((3,0)\) 满足 \(9/9+0/16=1\)，是椭圆 \(x^{2}/9+y^{2}/16=1\) 的右顶点，在椭圆上；过椭圆上一点有且仅有一条切线（此处为 \(x=3\)），故对．\par\noindent (4) 对．椭圆是封闭曲线，直线与它联立消元后二次项系数恒不为 0（不存在像抛物线平行对称轴、双曲线平行渐近线那样的降次情形），故「恰一个公共点」\(\iff\)「方程有重根」\(\iff\)「相切」，与变式 2 中抛物线的反例恰成对照．\par\noindent 故四判断依次为「错、错、对、对」．}
\end{ansblock}}

\xiaojie{①识别：凡联立消元判位置、求公共点，以及「切」与「一个公共点」的辨析题，判归本题型.}
\bindp {\kaishu ②操作：联立消元后先看二次项系数（平行渐近线、平行对称轴时降次），再看 \(\Delta\)；相切即二次方程有重根.}
\bindp {\kaishu ③收束：判完位置须落到公共点坐标或个数；举反例要同时验「恰一公共点」与「切线方程不同」两件.}

% ---------- 探究点二 公共点个数分类与特位 ----------
\tjdnr{二}{公共点个数分类与特位}{例\textbf{4}}{中档(知识点二)}{已知双曲线 \(C\)：\(2x^{2}-y^{2}=2\) 与点 \(P(1,2)\)，讨论过点 \(P\) 的直线 \(l\) 的斜率的情况，使 \(l\) 与双曲线 \(C\) 分别有一个公共点、两个公共点、没有公共点．}
\xiexwei{14mm}

{\ansblockgrayfalse
\begin{ansblock}[2章-练-课时17-11]
% ans:2章-练-课时17-11
\ansitem{6}{一个公共点：k∈\{√2,−√2,3/2\} 或 l 的斜率不存在；两个公共点：k∈(−∞,−√2)∪(−√2,√2)∪(√2,3/2)；没有公共点：k∈(3/2,+∞)}
\ansnote{详解}{① \(l\) 垂直 \(x\) 轴（斜率不存在）：直线 \(x=1\) 代入 \(2x^{2}-y^{2}=2\) 得 \(y=0\)，唯一公共点 \((1,0)\)（该处切线恰为 \(x=1\)），有一个公共点．\par\noindent ② \(l\) 不垂直 \(x\) 轴：设 \(y-2=k(x-1)\)，代入 \(2x^{2}-y^{2}=2\) 并整理得 \((2-k^{2})x^{2}+2(k^{2}-2k)x-k^{2}+4k-6=0\)．\par\noindent (a) \(2-k^{2}=0\)，即 \(k=\pm\sqrt{2}\)（\(l\) 平行渐近线 \(y=\pm\sqrt{2}x\)）：方程降为一次 \(2(k^{2}-2k)x-k^{2}+4k-6=0\)，一次项系数 \(2(2-2k)=4(1-k)\neq 0\)（\(k=\pm\sqrt{2}\) 时），故恰有一个公共点．\par\noindent (b) \(2-k^{2}\neq 0\)：\(\Delta=4(k^{2}-2k)^{2}+4(2-k^{2})(k^{2}-4k+6)=4(12-8k)=48-32k\)．\(\Delta=0\) \(\iff k=3/2\)，一个公共点（相切）；\(\Delta>0\) \(\iff k<3/2\)，两个公共点（本情形下须剔除 \(k=\pm\sqrt{2}\)，已归 (a)）；\(\Delta<0\) \(\iff k>3/2\)，没有公共点．\par\noindent 综上：有一个公共点 \(\iff k\in\{\sqrt{2},-\sqrt{2},3/2\}\) 或斜率不存在；有两个公共点 \(\iff k\in(-\infty,-\sqrt{2})\cup(-\sqrt{2},\sqrt{2})\cup(\sqrt{2},3/2)\)；没有公共点 \(\iff k\in(3/2,+\infty)\)．（注意 \(3/2\) 不在 \(k=\pm\sqrt{2}\) 之内（\(\sqrt{2}\approx 1.414<1.5\)），分段无冲突．）}
\end{ansblock}}

\liB{变式\textbf{3}}{中档(知识点二)}{}{已知直线 \(l\)：\(y=kx+2\) 和椭圆 \(C\)：\(x^{2}/3+y^{2}/2=1\)．分别求直线 \(l\) 与椭圆 \(C\) 有两个公共点、只有一个公共点和没有公共点时 \(k\) 的取值范围．}
\xiexwei{9mm}

{\ansblockgrayfalse
\begin{ansblock}[2章-练-课时17-13]
% ans:2章-练-课时17-13
\ansitem{7}{两个公共点：k<−√6/3 或 k>√6/3；只有一个公共点：k=±√6/3；没有公共点：−√6/3<k<√6/3}
\ansnote{详解}{把 \(y=kx+2\) 代入 \(x^{2}/3+y^{2}/2=1\)，乘 6 得 \(2x^{2}+3(kx+2)^{2}=6\)，即 \((2+3k^{2})x^{2}+12kx+6=0\)．二次项系数 \(2+3k^{2}\geqslant 2>0\) 恒成立（椭圆无降次情形），故公共点个数完全由 \(\Delta\) 定：\(\Delta=(12k)^{2}-4(2+3k^{2})\times 6=144k^{2}-48-72k^{2}=72k^{2}-48\)．\par\noindent \(\Delta>0\) \(\iff k^{2}>48/72=2/3\) \(\iff|k|>\sqrt{6}/3\)，两个公共点；\(\Delta=0\) \(\iff k=\pm\sqrt{6}/3\)，只有一个公共点（相切）；\(\Delta<0\) \(\iff k^{2}<2/3\) \(\iff|k|<\sqrt{6}/3\)，没有公共点．（边界核对：直线恒过 \((0,2)\)，而 \((0,2)\) 在椭圆外（\(0+4/2=2>1\)），故存在相离与相切位置、与读数自洽．）}
\end{ansblock}}

\tjdnr{二}{公共点个数分类与特位}{例\textbf{5}}{中档(知识点二)}{设直线 \(l\)：\(y=kx+1\)，抛物线 \(C\)：\(y^{2}=4x\)，当 \(k\) 为何值时，\(l\) 与 \(C\) 相切？相交？相离？}
\xiexwei{12mm}

{\ansblockgrayfalse
\begin{ansblock}[2章-练-课时17-12]
% ans:2章-练-课时17-12
\ansitem{8}{k=1 时相切；k<1 时相交；k>1 时相离}
\ansnote{详解}{联立 \(y=kx+1\) 与 \(y^{2}=4x\) 消 \(y\)：\((kx+1)^{2}=4x\)，即 \(k^{2}x^{2}+(2k-4)x+1=0\)．\par\noindent ① \(k\neq 0\)：二次项系数 \(k^{2}>0\)，为一元二次方程，\(\Delta=(2k-4)^{2}-4k^{2}=4k^{2}-16k+16-4k^{2}=16(1-k)\)．\(\Delta=0\) \(\iff k=1\)，相切；\(\Delta>0\) \(\iff k<1\) 且 \(k\neq 0\)，相交（两交点）；\(\Delta<0\) \(\iff k>1\)，相离．\par\noindent ② \(k=0\)：直线为 \(y=1\)，平行于对称轴 \(x\) 轴，代入得 \(1=4x\) 即 \(x=1/4\)，恰一个公共点 \((1/4,1)\)，属相交（不切）——此即例 2、变式 2 的退化特位，并入「相交」档．\par\noindent 综上：\(k=1\) 相切；\(k<1\)（含 \(k=0\)）相交；\(k>1\) 相离．}
\end{ansblock}}

\liB{变式\textbf{4}}{中档(知识点二)}{}{过原点的直线 \(l\) 与双曲线 \(x^{2}/4-y^{2}/3=1\) 相交于两点，求 \(l\) 的斜率的取值范围．}
\xiexwei{9mm}

{\ansblockgrayfalse
\begin{ansblock}[2章-练-课时17-14]
% ans:2章-练-课时17-14
\ansitem{9}{|k|<√3/2（即 k∈(−√3/2, √3/2)）}
\ansnote{详解}{设 \(l\)：\(y=kx\)（过原点的直线中 \(x=0\) 即 \(y\) 轴与双曲线无公共点，故可只议斜率存在者）．代入 \(x^{2}/4-y^{2}/3=1\)，乘 12 得 \(3x^{2}-4k^{2}x^{2}=12\)，即 \((3-4k^{2})x^{2}=12\)．\par\noindent ① \(3-4k^{2}>0\)（\(|k|<\sqrt{3}/2\)）：\(x^{2}=12/(3-4k^{2})>0\)，解得 \(x=\pm\sqrt{12/(3-4k^{2})}\)，对应两个不同交点（关于原点对称），符合「相交于两点」．\par\noindent ② \(3-4k^{2}\leqslant 0\)（\(|k|\geqslant\sqrt{3}/2\)）：当 \(=0\)（\(l\) 与渐近线 \(y=\pm(\sqrt{3}/2)x\) 平行，且过原点即为渐近线本身）时方程为 \(0=12\) 无解；当 \(<0\) 时 \(x^{2}<0\) 无实数解；两种都无公共点．\par\noindent 故所求范围为 \(|k|<\sqrt{3}/2\)．（几何直观：过原点（双曲线中心）的直线只要比两条渐近线陡即与双曲线交于两支上对称两点；渐近线本身与曲线零公共点，正是例 4 中 \(k=\pm\sqrt{2}\) 型降次的对照位．）}
\end{ansblock}}

\tjdnr{二}{公共点个数分类与特位}{例\textbf{6}}{简单(知识点二)}{过双曲线 \(x^{2}/20-y^{2}/5=1\) 的右焦点的直线被双曲线所截得的弦长为 \(\sqrt{5}\)，这样的直线有几条？}
\xiexwei{9mm}

{\ansblockgrayfalse
\begin{ansblock}[2章-练-课时17-10]
% ans:2章-练-课时17-10
\ansitem{10}{1}
\ansnote{详解}{\(a^{2}=20\)、\(b^{2}=5\)，\(a=2\sqrt{5}\)、\(c^{2}=25\) 即右焦点 \(F(5,0)\)．过 \(F\) 的直线（除 \(x\) 轴外）可设为 \(x=my+5\)，代入 \(x^{2}-4y^{2}=20\) 得 \((my+5)^{2}-4y^{2}=20\)，整理为 \((m^{2}-4)y^{2}+10my+5=0\)．\par\noindent ① \(m^{2}=4\)（\(|m|=2\)，即直线斜率 \(|k|=1/2=b/a\)，平行渐近线）：二次项消失，化为一次方程 \(10my+5=0\)，只有一个公共点，不构成弦．\par\noindent ② \(m^{2}<4\)（斜率 \(|k|>1/2\)，较渐近线陡，两点同在右支）：\(y_{1}+y_{2}=-10m/(m^{2}-4)\)、\(y_{1}y_{2}=5/(m^{2}-4)\)，故 \((y_{1}-y_{2})^{2}=[100m^{2}-20(m^{2}-4)]/(m^{2}-4)^{2}=80(m^{2}+1)/(m^{2}-4)^{2}\)，\(|AB|=\sqrt{1+m^{2}}\,|y_{1}-y_{2}|=4\sqrt{5}(m^{2}+1)/(4-m^{2})\)．记 \(u=m^{2}\in[0,4)\)，\(f(u)=-1+5/(4-u)\) 随 \(u\) 单调增，即 \(|AB|\) 随 \(u\) 单调增，最小在 \(u=0\)，值 \(4\sqrt{5}\times(1/4)=\sqrt{5}\)，且仅此一处；\(u=0\) 即直线 \(x=5\)（垂直 \(x\) 轴的通径），其长 \(2b^{2}/a=2\times 5/(2\sqrt{5})=\sqrt{5}\) \ding{51}．\par\noindent ③ \(m^{2}>4\)（斜率 \(|k|<1/2\)，较渐近线平，两点分居两支）：\(|AB|=4\sqrt{5}(u+1)/(u-4)\)，而 \((u+1)/(u-4)=1+5/(u-4)\) 随 \(u\) 单调减，值域为 \((4\sqrt{5},+\infty)\)（\(u\to 4^{+}\) 时趋于 \(+\infty\)，\(u\to+\infty\) 时趋于 \(4\sqrt{5}\)＝实轴长但取不到，水平线 \(y=0\) 单列如下），恒大于 \(4\sqrt{5}>\sqrt{5}\)，取不到 \(\sqrt{5}\)．另单独核直线 \(y=0\)（\(x\) 轴，未被 \(x=my+5\) 覆盖）：它与双曲线交于两顶点，弦长 \(2a=4\sqrt{5}\neq\sqrt{5}\)．\par\noindent 综上，被截弦长等于 \(\sqrt{5}\) 的直线只有通径 \(x=5\) 这一条，故填 1．}
\end{ansblock}}

\xiaojie{①识别：凡含参直线与曲线公共点个数讨论、弦长计数、过顶点弦与渐近线界的问题，判归本题型.}
\bindp {\kaishu ②操作：斜率不存在单独议；二次项系数含参先讨论降次（平行渐近线），再对二次情形用 \(\Delta\)；右支（范围）型补积与和的符号条件.}
\bindp {\kaishu ③收束：结论按「一个／两个／没有」分段作答，区间开闭用端点特位（切、降次、顶点）核验.}

% ---------- 探究点三 弦长、焦点弦与弦中点 ----------
\tjdnr{三}{弦长、焦点弦与弦中点}{例\textbf{7}}{简单(知识点三)}{已知直线 \(x-2y+2=0\) 与椭圆 \(x^{2}+4y^{2}=4\) 相交于 \(A\)，\(B\) 两点，求线段 \(AB\) 的长．}
\xiexwei{9mm}

{\ansblockgrayfalse
\begin{ansblock}[2章-练-课时17-03]
% ans:2章-练-课时17-03
\ansitem{11}{√5}
\ansnote{详解}{直线给出 \(x=2y-2\)，代入 \(x^{2}+4y^{2}=4\)：\((2y-2)^{2}+4y^{2}=4\to 4y^{2}-8y+4+4y^{2}=4\to 8y^{2}-8y=0\to 8y(y-1)=0\)，得 \(y=0\) 或 \(y=1\)，对应 \(x=-2\) 或 \(x=0\)．故 \(A(-2,0)\)、\(B(0,1)\)，\(|AB|=\sqrt{(0+2)^{2}+(1-0)^{2}}=\sqrt{5}\)．（此法直接解出两交点，无需弦长公式；若用弦长公式，\(x=my-2\) 型联立给 \(|y_{1}-y_{2}|=1\)、\(\sqrt{1+m^{2}}=\sqrt{5}\) 之倒数形式，结果同为 \(\sqrt{5}\)．）}
\end{ansblock}}

\liB{变式\textbf{5}}{简单(知识点三)}{}{过双曲线 \(x^{2}/3-y^{2}/6=1\) 的右焦点作倾斜角为 \(30^{\circ}\) 的直线 \(l\)，直线 \(l\) 与双曲线交于不同的两点 \(A\)，\(B\)，则 \(AB\) 的长为\kongbai{}．}
\xiexwei{9mm}

{\ansblockgrayfalse
\begin{ansblock}[2章-练-课时17-09]
% ans:2章-练-课时17-09
\ansitem{12}{16√3/5}
\ansnote{详解}{\(a^{2}=3\)、\(b^{2}=6\)，\(c^{2}=a^{2}+b^{2}=9\)，右焦点 \(F_{2}(3,0)\)．倾斜角 \(30^{\circ}\) 得 \(k=\tan 30^{\circ}=\sqrt{3}/3\)，直线 \(l\)：\(y=(\sqrt{3}/3)(x-3)\)．代入 \(x^{2}/3-y^{2}/6=1\)：以 \(y^{2}=(x-3)^{2}/3\) 代，得 \(x^{2}/3-(x-3)^{2}/18=1\)，乘 18 整理为 \(6x^{2}-(x-3)^{2}=18\)，即 \(5x^{2}+6x-27=0\)（二次项系数 \(5\neq 0\)，故 \(l\) 不与渐近线平行；\(\Delta=36+540=576>0\)，确两交点）．韦达：\(x_{1}+x_{2}=-6/5\)、\(x_{1}x_{2}=-27/5\)，故 \((x_{1}-x_{2})^{2}=(x_{1}+x_{2})^{2}-4x_{1}x_{2}=36/25+108/5=576/25\)，\(|x_{1}-x_{2}|=24/5\)．弦长 \(|AB|=\sqrt{1+k^{2}}\,|x_{1}-x_{2}|=\sqrt{1+1/3}\times 24/5=(2/\sqrt{3})\times(24/5)=48/(5\sqrt{3})=16\sqrt{3}/5\)．}
\end{ansblock}}

\tjdnr{三}{弦长、焦点弦与弦中点}{例\textbf{8}}{简单(知识点三)}{已知直线 \(AB\) 过抛物线 \(y^{2}=4x\) 的焦点，交抛物线于 \(A(x_{1},y_{1})\)，\(B(x_{2},y_{2})\) 两点，若 \(x_{1}+x_{2}=5\)，则 \(|AB|\) 等于\kongbai{}．}
\xiexwei{9mm}

\begin{ansblock}[2章-练-课时17-06]
% ans:2章-练-课时17-06
\ansitem{13}{7}
\ansnote{详解}{\(2p=4\)、\(p=2\)，焦点 \(F(1,0)\)、准线 \(x=-1\)．由抛物线定义，焦半径等于点到准线的距离：\(|AF|=x_{1}+1\)、\(|BF|=x_{2}+1\)．因 \(A\)、\(F\)、\(B\) 共线且 \(F\) 在 \(A\)、\(B\) 之间（焦点弦的两端点分居 \(x\) 轴两侧），\(|AB|=|AF|+|BF|=x_{1}+x_{2}+2=5+2=7\)．（即焦点弦公式 \(|AB|=x_{1}+x_{2}+p=5+2=7\)；按定义补出焦半径一步，防公式误记 \(p/2\)．）}
\end{ansblock}

\liB{变式\textbf{6}}{难(知识点三)}{}{过抛物线 \(x^{2}=2py\)（\(p>0\)）的焦点 \(F\) 作倾角为 \(30^{\circ}\) 的直线，与抛物线分别交于 \(A\)、\(B\) 两点（\(A\) 在 \(y\) 轴左侧），则 \(|AF|/|FB|=\)\kongbai{}．}
\xiexwei{12mm}

{\ansblockgrayfalse
\begin{ansblock}[2章-练-课时17-16]
% ans:2章-练-课时17-16
\ansitem{14}{1/3}
\ansnote{详解}{焦点 \(F(0,p/2)\)、准线 \(y=-p/2\)．倾角 \(30^{\circ}\) 即直线斜率 \(\tan 30^{\circ}=\sqrt{3}/3\)，直线 \(l\)：\(y=(\sqrt{3}/3)x+p/2\)，等价地 \(x=\sqrt{3}(y-p/2)\)．代入 \(x^{2}=2py\) 得 \(3(y-p/2)^{2}=2py\)，展开：\(3y^{2}-3py+3p^{2}/4-2py=0\)，即 \(3y^{2}-5py+3p^{2}/4=0\)，除以 3：\(y^{2}-(5p/3)y+p^{2}/4=0\)（\(\Delta=25p^{2}/9-p^{2}=16p^{2}/9\)，两根 \(y=(5p/3\pm 4p/3)/2\)，即 \(p/6\) 与 \(3p/2\)）．\par\noindent 为免 \(y\) 轴左右两侧混判，改以 \(x\) 为主元：把 \(y=x/\sqrt{3}+p/2\) 代 \(x^{2}=2py\) 得 \(x^{2}=2p(x/\sqrt{3}+p/2)=(2p/\sqrt{3})x+p^{2}\)，即 \(x^{2}-(2p/\sqrt{3})x-p^{2}=0\)．解之：\(x=(p/\sqrt{3})\pm(1/2)\sqrt{16p^{2}/3}=(p/\sqrt{3})\pm(2p/\sqrt{3})\)，故 \(x=3p/\sqrt{3}=\sqrt{3}p\) 或 \(x=-p/\sqrt{3}\)．\par\noindent 定位辨析：题设「\(A\) 在 \(y\) 轴左侧」即 \(x_{A}<0\)，故只能取负根 \(x_{A}=-p/\sqrt{3}\)（此时 \(y_{A}=x_{A}^{2}/(2p)=(p^{2}/3)/(2p)=p/6\)），而 \(x_{B}=\sqrt{3}p\)（\(y_{B}=3p^{2}/(2p)=3p/2\)）——此两步不可互换，否则比值倒置为 3．\par\noindent 焦半径（开口向上：\(|MF|=y_{M}+p/2\)，到准线的距离，见课前预习）：\(|AF|=y_{A}+p/2=p/6+p/2=2p/3\)；\(|FB|=y_{B}+p/2=3p/2+p/2=2p\)．\par\noindent 故 \(|AF|/|FB|=(2p/3)/(2p)=1/3\)．（互核一（代回曲线）：\(A(-p/\sqrt{3},p/6)\) 使 \(x^{2}=p^{2}/3=2p\times p/6\) \ding{51}；\(B(\sqrt{3}p,3p/2)\) 使 \(x^{2}=3p^{2}=2p\times(3p/2)\) \ding{51}．互核二（焦点弦坐标积）：本题抛物线开口向上，焦点弦之对偶形式为 \(x_{1}x_{2}=-p^{2}\)，本处 \((-p/\sqrt{3})\times\sqrt{3}p=-p^{2}\) \ding{51}．互核三（弦长）：\(|AF|+|FB|=2p/3+2p=8p/3\)，而 \(|AB|=y_{A}+y_{B}+p=p/6+3p/2+p=8p/3\)，与两点距 \(\sqrt{(4p/\sqrt{3})^{2}+(4p/3)^{2}}=\sqrt{64p^{2}/9}=8p/3\) 合 \ding{51}．）}
\end{ansblock}}

\tjdnr{三}{弦长、焦点弦与弦中点}{例\textbf{9}}{简单(知识点三)}{直线 \(y=x+1\) 被椭圆 \(x^{2}/4+y^{2}/2=1\) 所截得线段的中点的坐标是\nobreak(\kongwei)}

\bindopt {\fontsize{10.5pt}{\qpTJgLeadLiti}\selectfont \makebox[\dimexpr0.5\linewidth-1em\relax][l]{A．\((2/3,5/3)\)}\makebox[\dimexpr0.5\linewidth-1em\relax][l]{B．\((4/3,7/3)\)}\par}

\bindopt {\fontsize{10.5pt}{\qpTJgLeadLiti}\selectfont \makebox[\dimexpr0.5\linewidth-1em\relax][l]{C．\((-2/3,1/3)\)}\makebox[\dimexpr0.5\linewidth-1em\relax][l]{D．\((-13/2,-17/2)\)}\par}

{\ansblockgrayfalse
\begin{ansblock}[2章-练-课时17-08]
% ans:2章-练-课时17-08
\ansitem{15}{C}
\ansnote{详解}{联立消 \(y\)：\(x^{2}/4+(x+1)^{2}/2=1\)，乘 4 得 \(x^{2}+2(x+1)^{2}=4\)，即 \(3x^{2}+4x-2=0\)（\(\Delta=16-4\times 3\times(-2)=40>0\)，故确相交于两点）．设 \(A(x_{1},y_{1})\)、\(B(x_{2},y_{2})\)、中点 \(M(x_{0},y_{0})\)，韦达得 \(x_{1}+x_{2}=-4/3\)，故 \(x_{0}=(x_{1}+x_{2})/2=-2/3\)；中点在直线上，\(y_{0}=x_{0}+1=1/3\)．所以中点 \((-2/3,1/3)\)，选 C．（本席为点差法综合的铺垫位，正文按常规联立韦达书写．）}
\end{ansblock}}

\liB{变式\textbf{7}}{难(知识点三)}{}{已知椭圆 \(x^{2}/6+y^{2}/3=1\)，直线 \(x-y+1=0\) 与该椭圆交于 \(A\)，\(B\) 两点，求线段 \(AB\) 中点的坐标和线段 \(AB\) 的长．}
\xiexwei{12mm}

{\ansblockgrayfalse
\begin{ansblock}[2章-练-课时17-15]
% ans:2章-练-课时17-15
\ansitem{16}{中点 (−2/3, 1/3)；|AB|=8√2/3}
\ansnote{详解}{直线写成 \(y=x+1\)，代入 \(x^{2}/6+y^{2}/3=1\)，乘 6 得 \(x^{2}+2(x+1)^{2}=6\)，即 \(3x^{2}+4x-4=0\)（\(\Delta=16+48=64>0\)，确两交点）．设 \(A(x_{1},y_{1})\)、\(B(x_{2},y_{2})\)，则 \(x_{1}+x_{2}=-4/3\)、\(x_{1}x_{2}=-4/3\)．\par\noindent 中点：\(x_{0}=(x_{1}+x_{2})/2=-2/3\)，\(y_{0}=x_{0}+1=1/3\)，即中点 \((-2/3,1/3)\)．\par\noindent 弦长：\((x_{1}-x_{2})^{2}=(x_{1}+x_{2})^{2}-4x_{1}x_{2}=16/9+16/3=64/9\)，\(|x_{1}-x_{2}|=8/3\)；直线斜率 \(k=1\)，故 \(|AB|=\sqrt{1+k^{2}}\,|x_{1}-x_{2}|=\sqrt{2}\times 8/3=8\sqrt{2}/3\)．\par\noindent（回代验：方程 \(3x^{2}+4x-4=0\) 之根 \(x=2/3\)、\(x=-2\)，对应点 \((2/3,5/3)\)、\((-2,-1)\)：中点 \(((2/3-2)/2,(5/3-1)/2)=(-2/3,1/3)\) \ding{51}，\(|AB|=\sqrt{(2/3+2)^{2}+(5/3+1)^{2}}=(8/3)\sqrt{2}\) \ding{51}．）}
\end{ansblock}}

\xiaojie{①识别：凡给弦长、求弦长、弦中点与焦点弦比例的问题，判归本题型.}
\bindp {\kaishu ②操作：弦长 \(\sqrt{1+k^{2}}\,|x_{1}-x_{2}|\)（或以 \(y\) 为主元）；焦点弦用焦半径化坐标和；中点用韦达和之半＋中点在直线上.}
\bindp {\kaishu ③收束：结果回代曲线或用焦点弦结论（\(y_{1}y_{2}=-p^{2}\) 型）互核；含开口方向的题先定位再取根.}

% ---------- 探究点四 拓展·位置关系与公共点计数 ----------
\tjdnr{四}{拓展·位置关系与公共点计数}{例\textbf{10}}{简单(知识点二)}{如果直线 \(y=kx-1\) 与双曲线 \(x^{2}-y^{2}=4\) 没有公共点，求 \(k\) 的取值范围．}
\xiexwei{9mm}

{\ansblockgrayfalse
\begin{ansblock}[2章-拓-课时17-03]
% ans:2章-拓-课时17-03
\ansitem{17}{k<−√5/2 或 k>√5/2（|k|>√5/2）}
\ansnote{详解}{代入 \(x^{2}-(kx-1)^{2}=4\)，展开 \(x^{2}-k^{2}x^{2}+2kx-1=4\)，整理为 \((1-k^{2})x^{2}+2kx-5=0\)．\par\noindent ① \(1-k^{2}=0\)（\(k=\pm 1\)）：降为一次方程 \(2kx-5=0\)，有解 \(x=5/(2k)\)，即有公共点，不合「没有公共点」．\par\noindent ② \(1-k^{2}\neq 0\)：无公共点 \(\iff\Delta<0\)，\(\Delta=(2k)^{2}-4(1-k^{2})(-5)=4k^{2}+20(1-k^{2})=20-16k^{2}\)．令 \(20-16k^{2}<0\) 得 \(k^{2}>5/4\)，即 \(|k|>\sqrt{5}/2\)（此时 \(1-k^{2}\neq 0\) 自动满足，因 \(5/4>1\)）．\par\noindent 故 \(k\) 的取值范围为 \(k<-\sqrt{5}/2\) 或 \(k>\sqrt{5}/2\)．（几何校：双曲线 \(x^{2}/4-y^{2}/4=1\) 的渐近线斜率 \(\pm 1\)，直线恒过 \((0,-1)\)；\(|k|>\sqrt{5}/2\approx 1.118>1\) 属比渐近线陡、且与实轴交点在双曲线顶点内侧的情形，无交点自洽．）}
\end{ansblock}}

\liB{变式\textbf{8}}{中档(知识点二)}{}{过点 \((0,3)\) 作直线 \(l\)，它与双曲线 \(x^{2}/4-y^{2}/3=1\) 只有一个公共点，这样的直线有\nobreak(\kongwei)}

\bindopt {\fontsize{10.5pt}{\qpTJgLeadLiti}\selectfont \makebox[\dimexpr0.25\linewidth-0.5em\relax][l]{A．1条}\makebox[\dimexpr0.25\linewidth-0.5em\relax][l]{B．2条}\makebox[\dimexpr0.25\linewidth-0.5em\relax][l]{C．3条}\makebox[\dimexpr0.25\linewidth-0.5em\relax][l]{D．4条}\par}

{\ansblockgrayfalse
\begin{ansblock}[2章-拓-课时17-17]
% ans:2章-拓-课时17-17
\ansitem{18}{D（4条）}
\ansnote{详解}{① 斜率不存在：\(l\) 为 \(x=0\)（\(y\) 轴），代入得 \(-y^{2}/3=1\) 无实数解，公共点 0 个，不合．\par\noindent ② 斜率存在：设 \(l\)：\(y=kx+3\)，代入 \(x^{2}/4-y^{2}/3=1\)（乘 12）：\(3x^{2}-4(kx+3)^{2}=12\)，展开得 \((3-4k^{2})x^{2}-24kx-48=0\)．\par\noindent (a) \(3-4k^{2}\neq 0\)（\(l\) 不平行渐近线 \(y=\pm(\sqrt{3}/2)x\)）：一个公共点 \(\iff\Delta=0\)，\(\Delta=(-24k)^{2}+4(3-4k^{2})\times 48=576k^{2}+192(3-4k^{2})=576-192k^{2}\)，令其为 0 得 \(k^{2}=3\)，\(k=\pm\sqrt{3}\)——两条切线，符合．\par\noindent (b) \(3-4k^{2}=0\)，即 \(k=\pm\sqrt{3}/2\)：方程降为一次 \(-24kx-48=0\)，一次项系数 \(-24k\neq 0\)（\(k=\pm\sqrt{3}/2\) 时），有唯一解 \(x=-2/k\)，即 \(l\) 平行渐近线时与双曲线恰一个公共点——又两条，符合．\par\noindent 合计 \(2+2=4\) 条，选 D．（校：\(k=\pm\sqrt{3}\) 时 \(3-4k^{2}=3-12=-9\neq 0\) \ding{51} 两类不重；\(k=\pm\sqrt{3}/2\) 与 \(\pm\sqrt{3}\) 互异 \ding{51} 四类无重复．）}
\end{ansblock}}

\tjdnr{四}{拓展·位置关系与公共点计数}{例\textbf{11}}{中档(知识点二)}{已知直线 \(l\) 的方程为 \(y=kx-1\)，双曲线 \(C\) 的方程为 \(x^{2}-y^{2}=1\)．若直线 \(l\) 与双曲线 \(C\) 的右支相交于不同的两点，则实数 \(k\) 的取值范围是\nobreak(\kongwei)}

\bindopt {\fontsize{10.5pt}{\qpTJgLeadLiti}\selectfont \makebox[\dimexpr0.5\linewidth-1em\relax][l]{A．\((-\sqrt{2},\sqrt{2})\)}\makebox[\dimexpr0.5\linewidth-1em\relax][l]{B．\([1,\sqrt{2})\)}\par}

\bindopt {\fontsize{10.5pt}{\qpTJgLeadLiti}\selectfont \makebox[\dimexpr0.5\linewidth-1em\relax][l]{C．\([-\sqrt{2},\sqrt{2}]\)}\makebox[\dimexpr0.5\linewidth-1em\relax][l]{D．\((1,\sqrt{2})\)}\par}

{\ansblockgrayfalse
\begin{ansblock}[2章-拓-课时17-18]
% ans:2章-拓-课时17-18
\ansitem{19}{D}
\ansnote{详解}{代入 \(x^{2}-(kx-1)^{2}=1\)，展开 \(x^{2}-k^{2}x^{2}+2kx-1=1\)，整理为 \((1-k^{2})x^{2}+2kx-2=0\)．设两交点横坐标 \(x_{1}\)、\(x_{2}\)，「同在右支且为不同两点」须同时满足：\par\noindent (i) \(1-k^{2}\neq 0\)（保持二次，否则一次方程只一交点）；\par\noindent (ii) \(\Delta=(2k)^{2}-4(1-k^{2})(-2)=4k^{2}+8(1-k^{2})=8-4k^{2}>0\) \(\iff k^{2}<2\)；\par\noindent (iii) 两根之积 \(x_{1}x_{2}=-2/(1-k^{2})>0\) \(\iff 1-k^{2}<0\) \(\iff k^{2}>1\)（两根同号）；\par\noindent (iv) 两根之和 \(x_{1}+x_{2}=-2k/(1-k^{2})=2k/(k^{2}-1)>0\) \(\iff k>0\)（结合 (iii) 同为正，即在右支 \(x\geqslant 1\)）．\par\noindent 合 (ii)(iii)(iv)：\(1<k^{2}<2\) 且 \(k>0\)，得 \(1<k<\sqrt{2}\)，即 \((1,\sqrt{2})\)，选 D．（校：此范围内 \(1-k^{2}\neq 0\) 自动满足；端点 \(k=1\) 使二次项消失（一交点）、\(k=\sqrt{2}\) 使 \(\Delta=0\)（切，一交点），故两端均开，排除 B 之左闭．）}
\end{ansblock}}

\liB{变式\textbf{9}}{简单(知识点二)}{}{若曲线 \(2x=\sqrt{4+y^{2}}\) 与直线 \(y=m(x+1)\) 有公共点，则实数 \(m\) 的取值范围是\kongbai{}．}
\xiexwei{9mm}

{\ansblockgrayfalse
\begin{ansblock}[2章-拓-课时17-19]
% ans:2章-拓-课时17-19
\ansitem{20}{(−2,2)}
\ansnote{详解}{曲线侧：\(2x=\sqrt{4+y^{2}}\) 要求 \(x\geqslant 0\)，两边平方得 \(4x^{2}=4+y^{2}\)，即 \(x^{2}-y^{2}/4=1\) 且 \(x\geqslant 0\)（即该双曲线的右支，含右顶点 \((1,0)\)，渐近线 \(y=\pm 2x\)）．直线 \(y=m(x+1)\) 恒过定点 \((-1,0)\)．代入 \(4x^{2}-y^{2}=4\)：\(4x^{2}-m^{2}(x+1)^{2}=4\)，即 \((4-m^{2})x^{2}-2m^{2}x-(m^{2}+4)=0\)．\par\noindent ① \(|m|<2\)（\(4-m^{2}>0\)）：两根之积 \(=-(m^{2}+4)/(4-m^{2})<0\)，故有一正根；正根满足 \(x^{2}-y^{2}/4=1\) 且 \(x>0\)，即 \(x\geqslant 1\)，落在右支上，有公共点．\par\noindent ② \(|m|=2\)：方程降为 \(-2m^{2}x-(m^{2}+4)=0\)，解得 \(x=-(m^{2}+4)/(2m^{2})=-8/8=-1<0\)，不在右支，无公共点．\par\noindent ③ \(|m|>2\)（\(4-m^{2}<0\)）：积 \(>0\) 且和 \(=2m^{2}/(4-m^{2})<0\)，两根同为负，右支无公共点．\par\noindent 故有公共点 \(\iff|m|<2\)，即 \(m\in(-2,2)\)．}
\end{ansblock}}

\tjdnr{四}{拓展·位置关系与公共点计数}{例\textbf{12}}{中档(知识点一)}{已知斜率为 1 的直线 \(l\) 过椭圆 \(C\)：\(x^{2}/4+y^{2}=1\) 的右焦点，交椭圆于 \(A\)，\(B\) 两点，则弦 \(AB\) 的长为\nobreak(\kongwei)}

\lxopt{A．\(4/5\)}

\lxopt{B．\(6/5\)}

\lxopt{C．\(8/5\)}

\lxopt{D．\(13/5\)}

{\ansblockgrayfalse
\begin{ansblock}[2章-拓-课时17-01]
% ans:2章-拓-课时17-01
\ansitem{21}{C（弦长 8/5）}
\ansnote{详解}{\(a^{2}=4\)、\(b^{2}=1\)，\(c^{2}=a^{2}-b^{2}=3\)、\(c=\sqrt{3}\)，右焦点 \((\sqrt{3},0)\)．直线 \(l\)：\(y=x-\sqrt{3}\)．代入 \(x^{2}/4+y^{2}=1\)：\(x^{2}/4+(x-\sqrt{3})^{2}=1\)，乘 4 得 \(x^{2}+4(x^{2}-2\sqrt{3}x+3)=4\)，即 \(5x^{2}-8\sqrt{3}x+8=0\)．\(\Delta=(8\sqrt{3})^{2}-4\times 5\times 8=192-160=32>0\)，\(x_{1}+x_{2}=8\sqrt{3}/5\)、\(x_{1}x_{2}=8/5\)．\par\noindent 弦长 \(|AB|=\sqrt{1+k^{2}}\,|x_{1}-x_{2}|=\sqrt{2}\times\sqrt{\Delta}/5=\sqrt{2}\times\sqrt{32}/5=\sqrt{2}\times 4\sqrt{2}/5=8/5\)．（互核：\(|x_{1}-x_{2}|^{2}=(x_{1}+x_{2})^{2}-4x_{1}x_{2}=192/25-32/5=32/25\)，\(|AB|^{2}=2\times 32/25=64/25\)，\(|AB|=8/5\) \ding{51}．）故选 C．}
\end{ansblock}}

\xiaojie{①识别：凡无公共点求参、单公共点直线计数、交限定支（右支）求范围的问题，判归本题型.}
\bindp {\kaishu ②操作：降次（平行渐近线）单独成类；切线用 \(\Delta=0\)；右支（范围）型补积与和的符号（或 \(x\geqslant 1\)）条件.}
\bindp {\kaishu ③收束：条数题逐类穷举并核各类不重；范围题端点用降次位、切位、顶点位逐个核开闭.}

% ---------- 探究点五 拓展·焦点弦与弦长应用 ----------
\tjdnr{五}{拓展·焦点弦与弦长应用}{例\textbf{13}}{简单(知识点三)}{已知抛物线 \(y^{2}=4x\) 与过其焦点的斜率为 1 的直线交于 \(A\)，\(B\) 两点，\(O\) 为坐标原点，求 \(\overrightarrow{OA}\cdot\overrightarrow{OB}\)．}
\xiexwei{9mm}

{\ansblockgrayfalse
\begin{ansblock}[2章-拓-课时17-05]
% ans:2章-拓-课时17-05
\ansitem{22}{−3}
\ansnote{详解}{焦点 \(F(1,0)\)，直线 \(l\)：\(y=x-1\)．代入 \((x-1)^{2}=4x\) 得 \(x^{2}-2x+1=4x\)，即 \(x^{2}-6x+1=0\)（\(\Delta=36-4=32>0\)）．设 \(A(x_{1},y_{1})\)、\(B(x_{2},y_{2})\)，则 \(x_{1}+x_{2}=6\)、\(x_{1}x_{2}=1\)．又 \(y_{1}y_{2}=(x_{1}-1)(x_{2}-1)=x_{1}x_{2}-(x_{1}+x_{2})+1=1-6+1=-4\)．故 \(\overrightarrow{OA}\cdot\overrightarrow{OB}=x_{1}x_{2}+y_{1}y_{2}=1+(-4)=-3\)．（互核：以 \(y\) 为主元，\(x=y+1\) 代 \(y^{2}=4(y+1)\) 得 \(y^{2}-4y-4=0\)，\(y_{1}y_{2}=-4\) \ding{51}；\(x_{1}x_{2}=(y_{1}+1)(y_{2}+1)=-4+4+1=1\) \ding{51}．）}
\end{ansblock}}

\liB{变式\textbf{10}}{中档(知识点三)}{}{已知斜率为 2 的直线 \(AB\) 过抛物线 \(y^{2}=8x\) 的焦点，求弦 \(AB\) 的长．}
\xiexwei{9mm}

{\ansblockgrayfalse
\begin{ansblock}[2章-拓-课时17-07]
% ans:2章-拓-课时17-07
\ansitem{23}{10}
\ansnote{详解}{\(2p=8\)、\(p=4\)，焦点 \(F(2,0)\)、准线 \(x=-2\)．直线 \(l\)：\(y=2(x-2)\)．代入 \(y^{2}=8x\) 得 \(4(x-2)^{2}=8x\)，即 \((x-2)^{2}=2x\)，展开 \(x^{2}-4x+4=2x\)，得 \(x^{2}-6x+4=0\)（\(\Delta=36-16=20>0\)）．设 \(A(x_{1},y_{1})\)、\(B(x_{2},y_{2})\)，则 \(x_{1}+x_{2}=6\)．由抛物线定义（焦半径＝到准线距离，准线 \(x=-2\)）\(|AF|=x_{1}+2\)、\(|BF|=x_{2}+2\)，焦点弦 \(|AB|=|AF|+|BF|=x_{1}+x_{2}+4=6+4=10\)．（互核：以 \(y\) 为主元，\(x=y/2+2\) 代 \(y^{2}=8(y/2+2)\) 得 \(y^{2}-4y-16=0\)，\(y_{1}+y_{2}=4\)、\(y_{1}y_{2}=-16\)，\((y_{1}-y_{2})^{2}=16+64=80\)，\(|AB|=\sqrt{1+(1/2)^{2}}\,|y_{1}-y_{2}|=(\sqrt{5}/2)\times 4\sqrt{5}=10\) \ding{51}．）}
\end{ansblock}}

\tjdnr{五}{拓展·焦点弦与弦长应用}{例\textbf{14}}{中档(知识点三)}{过抛物线 \(y^{2}=8x\) 的焦点 \(F\) 的一条直线与此抛物线相交于 \(A\)，\(B\) 两点，已知 \(A(8,8)\)，求线段 \(AB\) 的中点到抛物线准线的距离．}
\xiexwei{9mm}

{\ansblockgrayfalse
\begin{ansblock}[2章-拓-课时17-12]
% ans:2章-拓-课时17-12
\ansitem{24}{25/4}
\ansnote{详解}{\(2p=8\)、\(p=4\)，焦点 \(F(2,0)\)，准线 \(x=-2\)．由 \(A(8,8)\)、\(F(2,0)\) 得直线 \(AB\) 的斜率 \(=(8-0)/(8-2)=4/3\)，方程 \(y=(4/3)(x-2)\)．代入 \(y^{2}=8x\)：\((16/9)(x-2)^{2}=8x\)，即 \(2(x-2)^{2}=9x\)，展开 \(2x^{2}-8x+8=9x\)，得 \(2x^{2}-17x+8=0\)（\(\Delta=289-64=225=15^{2}\)）．解得 \(x=(17\pm 15)/4\)，即 \(x=8\)（点 \(A\)）或 \(x=1/2\)，故 \(B(1/2,-2)\)（\(y=(4/3)(1/2-2)=-2\)）．\par\noindent 方法一（中点坐标）：中点横坐标 \(x_{0}=(8+1/2)/2=17/4\)，到准线 \(x=-2\) 的距离 \(=17/4+2=25/4\)．\par\noindent 方法二（焦半径/定义）：\(|AF|=x_{A}+p/2=8+2=10\)、\(|BF|=x_{B}+p/2=1/2+2=5/2\)，焦点弦上「中点到准线的距离＝\((|AF|+|BF|)/2=|AB|/2\)」（梯形中位线），得 \((10+5/2)/2=25/4\) \ding{51}．\par\noindent（校：\(y_{1}y_{2}=8\times(-2)=-16=-p^{2}\) \ding{51}，与焦点弦纵积定值相符．）}
\end{ansblock}}

\liB{变式\textbf{11}}{中档(知识点三)}{}{垂直于 \(x\) 轴的直线与抛物线 \(y^{2}=4x\) 交于 \(A\)，\(B\) 两点，且 \(|AB|=4\sqrt{3}\)，求直线 \(AB\) 的方程．}
\xiexwei{9mm}

{\ansblockgrayfalse
\begin{ansblock}[2章-拓-课时17-13]
% ans:2章-拓-课时17-13
\ansitem{25}{x=3}
\ansnote{详解}{设直线为 \(x=x_{0}\)（\(x_{0}\geqslant 0\)，否则与抛物线无交点）．代入 \(y^{2}=4x_{0}\)，得 \(y=\pm 2\sqrt{x_{0}}\)，两交点 \(A(x_{0},2\sqrt{x_{0}})\)、\(B(x_{0},-2\sqrt{x_{0}})\)（关于 \(x\) 轴对称，故该弦被 \(x\) 轴垂直平分）．\(|AB|=2\times 2\sqrt{x_{0}}=4\sqrt{x_{0}}=4\sqrt{3}\)，故 \(\sqrt{x_{0}}=\sqrt{3}\)、\(x_{0}=3\)．所以直线 \(AB\) 的方程为 \(x=3\)．（回验：\(y^{2}=12\)、\(y=\pm 2\sqrt{3}\)、\(|AB|=4\sqrt{3}\) \ding{51}；此即通径型垂直弦，长 \(4\sqrt{3}\) 大于通径 \(2p=4\)，点 \((3,\pm 2\sqrt{3})\) 在曲线上 \ding{51}．）}
\end{ansblock}}

\tjdnr{五}{拓展·焦点弦与弦长应用}{例\textbf{15}}{中档(知识点三)}{已知直线 \(l\)：\(y=x-3\) 与抛物线 \(C\)：\(x^{2}=-8y\) 相交于 \(A\)，\(B\) 两点，且 \(O\) 为坐标原点．}

\bindp (1) 求弦长 \(|AB|\) 以及线段 \(AB\) 的中点坐标；

\bindp (2) 判断 \(OA\perp OB\) 是否成立，并说明理由．
\xiexwei{14mm}

{\ansblockgrayfalse
\begin{ansblock}[2章-拓-课时17-08]
% ans:2章-拓-课时17-08
\ansitem{26}{(1) |AB|=8√5，中点 (−4,−7)；(2) 不成立（OA·OB=−15≠0）}
\ansnote{详解}{(1) 代入 \(x^{2}=-8(x-3)\)：\(x^{2}=-8x+24\)，即 \(x^{2}+8x-24=0\)（\(\Delta=64+96=160>0\)）．设 \(A(x_{1},y_{1})\)、\(B(x_{2},y_{2})\)，则 \(x_{1}+x_{2}=-8\)、\(x_{1}x_{2}=-24\)．弦长 \(|AB|=\sqrt{1+1^{2}}\,|x_{1}-x_{2}|=\sqrt{2}\times\sqrt{\Delta}=\sqrt{2}\times\sqrt{160}=\sqrt{320}=8\sqrt{5}\)．中点 \(x_{0}=(x_{1}+x_{2})/2=-4\)，\(y_{0}=x_{0}-3=-7\)，即 \((-4,-7)\)．（根之实值 \(x=-4\pm 2\sqrt{10}\)，对应 \(y=-7\pm 2\sqrt{10}\)，代 \(x^{2}=-8y\) 校：\((-4\pm 2\sqrt{10})^{2}=56\mp 16\sqrt{10}\)，\(-8(-7\pm 2\sqrt{10})=56\mp 16\sqrt{10}\) \ding{51}．）\par\noindent (2) \(y_{1}y_{2}=(x_{1}-3)(x_{2}-3)=x_{1}x_{2}-3(x_{1}+x_{2})+9=-24+24+9=9\)．故 \(\overrightarrow{OA}\cdot\overrightarrow{OB}=x_{1}x_{2}+y_{1}y_{2}=-24+9=-15\neq 0\)，所以 \(\overrightarrow{OA}\) 与 \(\overrightarrow{OB}\) 不垂直，「\(OA\perp OB\)」不成立．（若按斜率复核：\(k_{OA}\cdot k_{OB}=(y_{1}y_{2})/(x_{1}x_{2})=9/(-24)=-3/8\neq -1\)，同得否．）}
\end{ansblock}}

\liB{变式\textbf{12}}{中档(知识点三)}{}{已知斜率为 2 的直线 \(l\) 与抛物线 \(y^{2}=4x\) 相交于 \(A\)，\(B\) 两点，如果线段 \(AB\) 的长等于 5，求直线 \(l\) 的方程．}
\xiexwei{9mm}

{\ansblockgrayfalse
\begin{ansblock}[2章-拓-课时17-11]
% ans:2章-拓-课时17-11
\ansitem{27}{y=2x−2}
\ansnote{详解}{设 \(l\)：\(y=2x+m\)．因抛物线以 \(x=y^{2}/4\) 表出，改以 \(y\) 为主元：由 \(y=2x+m\) 得 \(x=(y-m)/2\)，代入 \(y^{2}=4x=2(y-m)\)，整理为 \(y^{2}-2y+2m=0\)（二次项系数 \(1\neq 0\)）．相交于两点须 \(\Delta=(-2)^{2}-4\times 2m=4-8m>0\)，即 \(m<1/2\)；此时 \(y_{1}+y_{2}=2\)、\(y_{1}y_{2}=2m\)，\((y_{1}-y_{2})^{2}=(y_{1}+y_{2})^{2}-4y_{1}y_{2}=4-8m\)．\par\noindent 弦长用 \(y\) 的差：\(|AB|=\sqrt{1+(1/k)^{2}}\,|y_{1}-y_{2}|=\sqrt{1+1/4}\times\sqrt{4-8m}=(\sqrt{5}/2)\sqrt{4-8m}\)．令其等于 5：\(\sqrt{4-8m}=10/\sqrt{5}=2\sqrt{5}\)，平方得 \(4-8m=20\)，故 \(m=-2\)（\(<1/2\) \ding{51} 合相交条件）．\par\noindent 所以直线 \(l\) 的方程为 \(y=2x-2\)．（回代以 \(x\) 为主元复核：\((2x-2)^{2}=4x\to 4x^{2}-8x+4=4x\to x^{2}-3x+1=0\)，\(|x_{1}-x_{2}|=\sqrt{9-4}=\sqrt{5}\)，\(|AB|=\sqrt{1+4}\times\sqrt{5}=5\) \ding{51}．）}
\end{ansblock}}

\tjdnr{五}{拓展·焦点弦与弦长应用}{例\textbf{16}}{简单(知识点三)}{求圆心在抛物线 \(y^{2}=2x\) 上且与 \(x\) 轴和该抛物线的准线都相切的圆的方程．}
\xiexwei{12mm}

{\ansblockgrayfalse
\begin{ansblock}[2章-拓-课时17-04]
% ans:2章-拓-课时17-04
\ansitem{28}{(x−1/2)²+(y−1)²=1 或 \hspace{0pt plus 1.5em}(x−1/2)²+(y+1)²=1（两圆）}
\ansnote{详解}{\(y^{2}=2x\) 中 \(2p=2\)、\(p=1\)，准线 \(l\)：\(x=-1/2\)．设圆心 \(C(a,b)\) 在抛物线上，则 \(b^{2}=2a\)（故 \(a\geqslant 0\)）、半径 \(r\)．与 \(x\) 轴相切 \(\Rightarrow r=|b|\)；与准线 \(x=-1/2\) 相切 \(\Rightarrow r=a+1/2\)（\(a\geqslant 0\) 故 \(a-(-1/2)=a+1/2\)）．两式相等：\(|b|=a+1/2\)，以 \(a=b^{2}/2\) 代得 \(|b|=b^{2}/2+1/2\)，即 \(|b|^{2}-2|b|+1=0\)，\((|b|-1)^{2}=0\)，故 \(|b|=1\)、\(b=\pm 1\)，进而 \(a=1/2\)、\(r=1\)．\par\noindent 所以圆有两解：\((x-1/2)^{2}+(y-1)^{2}=1\) 与 \((x-1/2)^{2}+(y+1)^{2}=1\)．（回验：圆心 \((1/2,1)\) 在抛物线上 \(1=2\times 1/2\) \ding{51}；到 \(x\) 轴距离 \(1=r\) \ding{51}；到准线距离 \(1/2+1/2=1=r\) \ding{51}．）}
\end{ansblock}}

\xiaojie{①识别：凡焦点弦长、弦中点到准线距离、垂直弦与相切圆等焦半径应用问题，判归本题型.}
\bindp {\kaishu ②操作：焦半径 \(|PF|=x_{P}+p/2\) 当第一工具；弦长公式选主元（与开口垂直的方向常更快）；相切圆抓住「圆心到线距离＝半径」.}
\bindp {\kaishu ③收束：双解（\(b=\pm\) 型）逐个回验三条件；弦长定值反求时平方前先核 Δ 与相交前提.}

% ---------- 探究点六 拓展·证明与综合收难 ----------
\tjdnr{六}{拓展·证明与综合收难}{例\textbf{17}}{中档(知识点三)}{过抛物线 \(y^{2}=2px\)（\(p>0\)）的焦点的一条直线与抛物线相交，且两个交点的纵坐标为 \(y_{1}\)，\(y_{2}\)，求证：\(y_{1}y_{2}=-p^{2}\)．}
\xiexwei{10.5mm}

{\ansblockgrayfalse
\begin{ansblock}[2章-拓-课时17-06]
% ans:2章-拓-课时17-06
\ansitem{29}{证明见详解}
\ansnote{详解}{焦点 \(F(p/2,0)\)．设过 \(F\) 的直线为 \(x=my+p/2\)：此写法已含垂直 \(x\) 轴的弦（\(m=0\)，即通径），而未含的只有 \(y=0\)（抛物线的对称轴），它与抛物线仅一个公共点，与题设「相交（有两个交点）」矛盾，故不损一般性．代入 \(y^{2}=2px\)：\(y^{2}=2p(my+p/2)=2pmy+p^{2}\)，整理为关于 \(y\) 的一元二次方程 \(y^{2}-2pmy-p^{2}=0\)．其二次项系数为 \(1\neq 0\)，且 \(\Delta=(2pm)^{2}+4p^{2}=4p^{2}(m^{2}+1)>0\)（\(p>0\)），故确两交点，与题设相符．由韦达定理，\(y_{1}y_{2}=-p^{2}/1=-p^{2}\)．（特位核验：\(m=0\) 通径时 \(y^{2}=p^{2}\)、\(y=\pm p\)、积 \(-p^{2}\) \ding{51}；\(m=1\)、\(p=2\) 时 \(y^{2}-4y-4=0\)、积 \(-4=-p^{2}\) \ding{51}，与例 13 之 \(y_{1}y_{2}=-4\) 同．）}
\end{ansblock}}

\liB{变式\textbf{13}}{中档(知识点二)}{}{已知直线 \(l\) 的斜率与双曲线 \(C\) 的渐近线的斜率相等，求证：直线 \(l\) 与双曲线 \(C\) 最多只有一个公共点．}
\xiexwei{12mm}

{\ansblockgrayfalse
\begin{ansblock}[2章-拓-课时17-09]
% ans:2章-拓-课时17-09
\ansitem{30}{证明见详解}
\ansnote{详解}{双曲线分两种标准形，证法同，先取 \(C\)：\(x^{2}/a^{2}-y^{2}/b^{2}=1\)（\(a>0,b>0\)），其渐近线为 \(y=\pm(b/a)x\)．题设「\(l\) 的斜率与渐近线的斜率相等」即 \(l\) 与某条渐近线平行或重合，不妨设 \(l\)：\(y=(b/a)x+m\)（\(m\) 为任意实数，\(m=0\) 时 \(l\) 就是渐近线本身）．代入 \(C\)：\(x^{2}/a^{2}-((b/a)x+m)^{2}/b^{2}=1\)，而 \(((b/a)x+m)^{2}/b^{2}=(x/a+m/b)^{2}=x^{2}/a^{2}+2mx/(ab)+m^{2}/b^{2}\)，故二次项 \(x^{2}/a^{2}\) 相消，得 \(-2mx/(ab)-m^{2}/b^{2}=1\)，即 \(-(2m/(ab))x=1+m^{2}/b^{2}\)．\par\noindent ① \(m\neq 0\)：一次项系数 \(-2m/(ab)\neq 0\)，方程有唯一解 \(x=-ab(1+m^{2}/b^{2})/(2m)\)，即 \(l\) 与 \(C\) 恰一个公共点．\par\noindent ② \(m=0\)：等式左端为 0、右端为 1，矛盾，无解，即渐近线与 \(C\) 没有公共点．\par\noindent 两种情形公共点均不超过一个，故「最多只有一个公共点」．焦点在 \(y\) 轴上时（\(C\)：\(y^{2}/a^{2}-x^{2}/b^{2}=1\)，渐近线 \(y=\pm(a/b)x\)）同法：设 \(y=(a/b)x+m\) 代入并以 \(x\) 为主元，二次项 \(x^{2}/b^{2}-x^{2}/b^{2}=0\) 相消，化为一次方程或矛盾式，结论同上．（本条与例 4、变式 4 之降次情形同一事理，可作其原理说明．）}
\end{ansblock}}

\tjdnr{六}{拓展·证明与综合收难}{例\textbf{18}}{中档(知识点二)}{求过点 \(A(0,p)\) 且与抛物线 \(y^{2}=2px\)（\(p>0\)）只有一个公共点的直线的方程．}
\xiexwei{14mm}

{\ansblockgrayfalse
\begin{ansblock}[2章-拓-课时17-10]
% ans:2章-拓-课时17-10
\ansitem{31}{三条：x=0；y=p；y=(1/2)x+p（即 x−2y+2p=0）}
\ansnote{详解}{分「斜率不存在」「斜率为 0」「斜率存在且非 0」三类穷举．\par\noindent ① 斜率不存在：直线 \(x=0\)（即 \(y\) 轴），代 \(y^{2}=2p\times 0=0\) 得 \(y=0\)，唯一公共点 \((0,0)\)（顶点），符合．\par\noindent ② 斜率为 0：直线 \(y=p\)（过 \(A(0,p)\) 且平行对称轴 \(x\) 轴），代 \(p^{2}=2px\) 得 \(x=p/2\)，唯一公共点 \((p/2,p)\)，符合——此即例 2／变式 2 的退化型「一公共点而不切」．\par\noindent ③ 斜率存在且非 0：设 \(y=kx+p\)（\(k\neq 0\)），代 \((kx+p)^{2}=2px\)，得 \(k^{2}x^{2}+(2kp-2p)x+p^{2}=0\)（二次项系数 \(k^{2}\neq 0\)）．只有一个公共点 \(\iff\Delta=0\)：\(\Delta=(2kp-2p)^{2}-4k^{2}p^{2}=4p^{2}[(k-1)^{2}-k^{2}]=4p^{2}(1-2k)=0\)，故 \(k=1/2\)，切线为 \(y=(1/2)x+p\)，唯一公共点为 \(x=(2p-2kp)/(2k^{2})=(2p-p)/(1/2)=2p\)、\(y=2p\)，即切点 \((2p,2p)\)（校：\((2p)^{2}=4p^{2}=2p\times 2p\) \ding{51}）．\par\noindent 综上满足条件的直线共三条：\(x=0\)、\(y=p\)、\(y=(1/2)x+p\)．（几何校：\(A(0,p)\) 在抛物线外（\(p^{2}>0=2p\times 0\)），故恰一条切线；另两条分别由「与对称轴平行」「过顶点的竖直割线」两个退化方向提供，条数与读数合．）}
\end{ansblock}}

\liB{变式\textbf{14}}{中档(知识点三)}{}{过抛物线的焦点的一条直线与它交于 \(P\)，\(Q\) 两点，过点 \(P\) 和此抛物线顶点的直线与抛物线的准线交于点 \(M\)，求证：直线 \(MQ\) 平行于此抛物线的对称轴．}
\xiexwei{12mm}

{\ansblockgrayfalse
\begin{ansblock}[2章-拓-课时17-14]
% ans:2章-拓-课时17-14
\ansitem{32}{证明见详解}
\ansnote{详解}{取抛物线标准形 \(y^{2}=2px\)（\(p>0\)），顶点 \(O(0,0)\)、焦点 \(F(p/2,0)\)、准线 \(x=-p/2\)、对称轴为 \(x\) 轴．设 \(P(x_{1},y_{1})\)、\(Q(x_{2},y_{2})\)，\(y_{1}y_{2}\neq 0\)（焦点弦的端点不为顶点，且由例 17 有 \(y_{1}y_{2}=-p^{2}\neq 0\)）．\par\noindent 直线 \(OP\) 的斜率 \(=y_{1}/x_{1}\)，而 \(y_{1}^{2}=2px_{1}\) 故 \(x_{1}=y_{1}^{2}/(2p)\)，斜率 \(=y_{1}/(y_{1}^{2}/(2p))=2p/y_{1}\)，即 \(OP\)：\(y=(2p/y_{1})x\)．\par\noindent 令 \(x=-p/2\)（准线），得 \(M\) 的纵坐标 \(y_{M}=(2p/y_{1})(-p/2)=-p^{2}/y_{1}\)．\par\noindent 由 \(y_{1}y_{2}=-p^{2}\) 得 \(y_{2}=-p^{2}/y_{1}\)，于是 \(y_{M}=y_{2}\)．\par\noindent \(M(-p/2,y_{M})\) 与 \(Q(x_{2},y_{2})\) 纵坐标相同，而横坐标不同（\(x_{2}\geqslant 0>-p/2\)），故直线 \(MQ\) 为水平线，即 \(MQ\) 平行于抛物线的对称轴 \(x\) 轴．（通径特位：\(PQ\perp x\) 轴时 \(x_{1}=x_{2}=p/2\)、\(y_{2}=-y_{1}\)，由 \(y_{1}^{2}=2p\times(p/2)=p^{2}\) 得 \(y_{M}=-p^{2}/y_{1}=-y_{1}=y_{2}\) \ding{51} 证法不变；焦点在 \(y\) 轴、开口向左等其他标准形同法（换元即可），结论为「\(MQ\) 平行于该抛物线的对称轴」．）}
\end{ansblock}}

\tjdnr{六}{拓展·证明与综合收难}{例\textbf{19}}{中档(知识点三)}{过抛物线的焦点 \(F\) 的一条直线与此抛物线相交于 \(P_{1}\)，\(P_{2}\) 两点．求证：以 \(P_{1}P_{2}\) 为直径的圆与该抛物线的准线相切．}
\xiexwei{12mm}

{\ansblockgrayfalse
\begin{ansblock}[2章-拓-课时17-15]
% ans:2章-拓-课时17-15
\ansitem{33}{证明见详解}
\ansnote{详解}{取标准形 \(y^{2}=2px\)（\(p>0\)），准线 \(l\)：\(x=-p/2\)．设 \(P_{1}(x_{1},y_{1})\)、\(P_{2}(x_{2},y_{2})\)，弦 \(P_{1}P_{2}\) 的中点为 \(M\)，则圆心为 \(M\)、半径 \(r=|P_{1}P_{2}|/2\)．\par\noindent 由抛物线定义，焦半径等于点到准线的距离：\(|P_{1}F|=x_{1}+p/2\)、\(|P_{2}F|=x_{2}+p/2\)．因 \(F\) 在线段 \(P_{1}P_{2}\) 上（焦点弦过焦点），\(|P_{1}P_{2}|=|P_{1}F|+|P_{2}F|=x_{1}+x_{2}+p\)，故 \(r=(x_{1}+x_{2}+p)/2\)．\par\noindent 又圆心 \(M\) 的横坐标为 \((x_{1}+x_{2})/2\)，\(M\) 到准线 \(x=-p/2\) 的距离 \(d=(x_{1}+x_{2})/2+p/2=(x_{1}+x_{2}+p)/2\)．\par\noindent 于是 \(d=r\)，即圆心到准线的距离恰等于半径，故以 \(P_{1}P_{2}\) 为直径的圆与准线相切．（另一表述：分别过 \(P_{1}\)、\(P_{2}\)、\(M\) 向准线作垂线，垂足为 \(H_{1}\)、\(H_{2}\)、\(H\)，则 \(|MH|\) 是直角梯形 \(P_{1}H_{1}H_{2}P_{2}\) 的中位线，\(|MH|=(|P_{1}F|+|P_{2}F|)/2=|P_{1}P_{2}|/2=r\)，同得相切——此即「中位线＝焦半径和之半」结构，例 14 方法二同族．）}
\end{ansblock}}

\liB{变式\textbf{15}}{中档(知识点三)}{}{设抛物线 \(C\)：\(y^{2}=4x\) 的焦点为 \(F\)，过 \(F\) 且斜率为 1 的直线 \(l\) 与 \(C\) 相交于 \(A\)，\(B\) 两点，求过点 \(A\)，\(B\) 且与 \(C\) 的准线相切的圆的方程．}
\xiexwei{14mm}

{\ansblockgrayfalse
\begin{ansblock}[2章-拓-课时17-16]
% ans:2章-拓-课时17-16
\ansitem{34}{(x−3)²+(y−2)²=16 或 (x−11)²+(y+6)²=144（两圆）}
\ansnote{详解}{焦点 \(F(1,0)\)、准线 \(x=-1\)，直线 \(l\)：\(y=x-1\)．代入 \((x-1)^{2}=4x\) 得 \(x^{2}-6x+1=0\)（\(\Delta=32>0\)），设 \(A(x_{1},y_{1})\)、\(B(x_{2},y_{2})\)，则 \(x_{1}+x_{2}=6\)、\(x_{1}x_{2}=1\)，\(x=3\pm 2\sqrt{2}\)．\par\noindent 先定必备数据：弦的中点 \(M=(3,2)\)；\(|AB|=\sqrt{2}\,|x_{1}-x_{2}|=\sqrt{2}\times\sqrt{36-4}=\sqrt{2}\times\sqrt{32}=8\)，故 \(|MA|=|MB|=4\)．\par\noindent 所求圆过 \(A\)、\(B\)，其圆心在 \(AB\) 的中垂线上：中垂线斜率 \(-1\)、过 \(M\)，即 \(y=-x+5\)，设圆心 \(C(t,5-t)\)．\par\noindent 圆的半径按「与准线 \(x=-1\) 相切」为 \(r=|t+1|\)；按「过 \(A\)」为 \(r^{2}=|CA|^{2}=|CM|^{2}+|MA|^{2}=2(t-3)^{2}+16\)（因 \(|CM|^{2}=(t-3)^{2}+(5-t-2)^{2}=2(t-3)^{2}\)）．\par\noindent 令两式相等：\((t+1)^{2}=2(t-3)^{2}+16\)，展开 \(t^{2}+2t+1=2t^{2}-12t+18+16\)，移项得 \(t^{2}-14t+33=0\)，\((t-3)(t-11)=0\)，故 \(t=3\) 或 \(t=11\)（均使 \(r=|t+1|>0\) \ding{51}，两解皆合）．\par\noindent \(t=3\)：圆心 \((3,2)=M\)、\(r=4\)，即 \((x-3)^{2}+(y-2)^{2}=16\)——此圆恰为以焦点弦 \(AB\) 为直径的圆，与例 19（以焦点弦为直径的圆与准线相切）相符．\par\noindent \(t=11\)：圆心 \((11,-6)\)、\(r=12\)，即 \((x-11)^{2}+(y+6)^{2}=144\)．\par\noindent 故满足条件的圆有两个：\((x-3)^{2}+(y-2)^{2}=16\) 与 \((x-11)^{2}+(y+6)^{2}=144\)．（验第二圆：取 \(A(3+2\sqrt{2},2+2\sqrt{2})\)，\(|CA|^{2}=(8-2\sqrt{2})^{2}+(-8-2\sqrt{2})^{2}=144=r^{2}\) \ding{51}；圆心到准线距离 \(11+1=12=r\) \ding{51}．）}
\end{ansblock}}

\tjdnr{六}{拓展·证明与综合收难}{例\textbf{20}}{中档(知识点二)}{已知双曲线 \(C\)：\(2x^{2}-y^{2}=2\)，过点 \(P(1,2)\) 的直线 \(l\) 与双曲线 \(C\) 交于 \(M\)、\(N\) 两点，若 \(P\) 为线段 \(MN\) 的中点，则弦长 \(|MN|\) 等于\nobreak(\kongwei)}

\bindopt {\fontsize{10.5pt}{\qpTJgLeadLiti}\selectfont \makebox[\dimexpr0.25\linewidth-0.5em\relax][l]{A．\(4\sqrt{2}/3\)}\makebox[\dimexpr0.25\linewidth-0.5em\relax][l]{B．\(3\sqrt{3}/4\)}\makebox[\dimexpr0.25\linewidth-0.5em\relax][l]{C．\(4\sqrt{3}\)}\makebox[\dimexpr0.25\linewidth-0.5em\relax][l]{D．\(4\sqrt{2}\)}\par}

{\ansblockgrayfalse
\begin{ansblock}[2章-拓-课时17-20]
% ans:2章-拓-课时17-20
\ansitem{35}{D}
\ansnote{详解}{若 \(l\) 斜率不存在（\(x=1\)）：代入得 \(2-y^{2}=2\)，\(y=0\)，只一交点，不合「交于两点且 \(P\) 为中点」；故设 \(l\)：\(y=k(x-1)+2\)．代入 \(2x^{2}-y^{2}=2\)：\((2-k^{2})x^{2}+2k(k-2)x-(k^{2}-4k+6)=0\)（其中一次项系数由 \(-2k(2-k)\) 化为 \(2k(k-2)\)、常数项为 \(-(2-k)^{2}-2=-(k^{2}-4k+6)\)）．须 \(2-k^{2}\neq 0\)．\par\noindent 中点条件：\(x_{1}+x_{2}=2\) \(\Rightarrow -2k(k-2)/(2-k^{2})=2\) \(\Rightarrow -2k^{2}+4k=4-2k^{2}\) \(\Rightarrow 4k=4\) \(\Rightarrow k=1\)（合 \(2-k^{2}=1\neq 0\) \ding{51}）．\par\noindent 此时方程为 \(x^{2}-2x-3=0\)，\(x_{1}+x_{2}=2\)、\(x_{1}x_{2}=-3\)，两根 \(3\)、\(-1\)，\(\Delta=4+12=16>0\) \ding{51} 两交点 \(M(3,4)\)、\(N(-1,0)\)（校：\(2\times 9-16=2\) \ding{51}；\(2\times 1-0=2\) \ding{51}）．\par\noindent \(|MN|=\sqrt{1+k^{2}}\,|x_{1}-x_{2}|=\sqrt{2}\times\sqrt{(x_{1}+x_{2})^{2}-4x_{1}x_{2}}=\sqrt{2}\times\sqrt{4+12}=\sqrt{2}\times 4=4\sqrt{2}\)，选 D．}
\end{ansblock}}

\liB{变式\textbf{16}}{简单(知识点二)}{}{直线 \(l\) 与双曲线 \(x^{2}-y^{2}/2=1\) 交于 \(A\)，\(B\) 两点，以 \(AB\) 为直径的圆 \(C\) 的方程为 \(x^{2}+y^{2}+2x+4y+m=0\)，则 \(m=\)\nobreak(\kongwei)}

\bindopt {\fontsize{10.5pt}{\qpTJgLeadLiti}\selectfont \makebox[\dimexpr0.25\linewidth-0.5em\relax][l]{A．\(-3\)}\makebox[\dimexpr0.25\linewidth-0.5em\relax][l]{B．\(3\)}\makebox[\dimexpr0.25\linewidth-0.5em\relax][l]{C．\(5-2\sqrt{2}\)}\makebox[\dimexpr0.25\linewidth-0.5em\relax][l]{D．\(2\sqrt{2}\)}\par}

{\ansblockgrayfalse
\begin{ansblock}[2章-拓-课时17-21]
% ans:2章-拓-课时17-21
\ansitem{36}{A（m=−3）}
\ansnote{详解}{圆 \(C\) 配方法：\((x+1)^{2}+(y+2)^{2}=5-m\)（须 \(m<5\)），圆心 \((-1,-2)\)、半径 \(r=\sqrt{5-m}\)．\(AB\) 是圆 \(C\) 的直径，故圆心即 \(AB\) 的中点：中点 \((-1,-2)\)．\par\noindent 求弦所在直线（按常规联立韦达，不引中点弦公式名目）：设 \(l\) 斜率为 \(k\)（若 \(l\) 为 \(x=-1\)，代入双曲线得 \(1-y^{2}/2=1\)、\(y=0\)，只一交点，不合），以 \(u=x+1\) 平移（则 \(x=u-1\)、\(y=ku-2\)，中点条件即 \(u_{1}+u_{2}=0\)）．代入 \(2x^{2}-y^{2}=2\)：\(2(u-1)^{2}-(ku-2)^{2}=2\to 2u^{2}-4u+2-k^{2}u^{2}+4ku-4=2\)，即 \((2-k^{2})u^{2}+(4k-4)u-4=0\)．由 \(u_{1}+u_{2}=0\) 得 \(-(4k-4)/(2-k^{2})=0\)，故 \(k=1\)（\(2-k^{2}=1\neq 0\) \ding{51}）．\par\noindent 于是 \(l\)：\(y=(x+1)-2=x-1\)，方程化为 \(u^{2}-4=0\)，\(u=\pm 2\to A(1,0)\)、\(B(-3,-4)\)（校：\(1-0=1\) \ding{51}；\(9-16/2=1\) \ding{51}），\(|AB|=\sqrt{4^{2}+4^{2}}=4\sqrt{2}\)．\par\noindent 由直径 \(|AB|^{2}=D^{2}+E^{2}-4F\)（此处 \(D=2\)、\(E=4\)、\(F=m\)）：\(|AB|^{2}=4+16-4m=20-4m\)，令 \((4\sqrt{2})^{2}=32=20-4m\)，解得 \(m=-3\)（\(<5\) \ding{51} 圆确存在）．故选 A．（互核：\(r=|AB|/2=2\sqrt{2}\)、\(r^{2}=8=5-m\Rightarrow m=-3\) \ding{51}；又 \(|CA|^{2}=(1+1)^{2}+(0+2)^{2}=8=r^{2}\) \ding{51}．）}
\end{ansblock}}

\tjdnr{六}{拓展·证明与综合收难}{例\textbf{21}}{中档(知识点三)}{已知直线 \(l\)：\(y=kx+1\) 与椭圆 \(x^{2}/2+y^{2}=1\) 交于 \(M\)，\(N\) 两点，且 \(|MN|=4\sqrt{2}/3\)，则 \(k=\)\kongbai{}．}
\xiexwei{9mm}

{\ansblockgrayfalse
\begin{ansblock}[2章-拓-课时17-02]
% ans:2章-拓-课时17-02
\ansitem{37}{±1}
\ansnote{详解}{联立消 \(y\)：\(x^{2}/2+(kx+1)^{2}=1\)，乘 2 得 \(x^{2}+2(k^{2}x^{2}+2kx+1)=2\)，即 \((1+2k^{2})x^{2}+4kx=0\)．二次项系数 \(1+2k^{2}>0\) 恒成立．设 \(M(x_{1},y_{1})\)、\(N(x_{2},y_{2})\)，则 \(x_{1}+x_{2}=-4k/(1+2k^{2})\)、\(x_{1}x_{2}=0\)（事实上直线过定点 \((0,1)\)，即椭圆上顶点，一交点恒为它）．\par\noindent 由 \(|MN|=4\sqrt{2}/3\) 得 \((x_{1}-x_{2})^{2}+(y_{1}-y_{2})^{2}=32/9\)，即 \((1+k^{2})(x_{1}-x_{2})^{2}=32/9\)；而 \((x_{1}-x_{2})^{2}=(x_{1}+x_{2})^{2}-4x_{1}x_{2}=16k^{2}/(1+2k^{2})^{2}\)，代入得 \((1+k^{2})\times 16k^{2}/(1+2k^{2})^{2}=32/9\)，约 16 并交叉相乘：\(9k^{2}(1+k^{2})=2(1+2k^{2})^{2}\)，展开 \(9k^{2}+9k^{4}=2+8k^{2}+8k^{4}\)，即 \(k^{4}+k^{2}-2=0\)，\((k^{2}+2)(k^{2}-1)=0\)，故 \(k^{2}=1\)、\(k=\pm 1\)．（回验 \(k=1\)：\((1+2)x^{2}+4x=0\to x=0\)、\(-4/3\)，\(|MN|=\sqrt{2}\times 4/3=4\sqrt{2}/3\) \ding{51}．）}
\end{ansblock}}

\xiaojie{①识别：凡焦点弦定值证明、渐近线降次原理、切线条数穷举与中点弦直径圆等证明综合题，判归本题型.}
\bindp {\kaishu ②操作：证明题设线不损一般性（补特位核验）；定值先算积（和）再约参；穷举按「斜率不存在／为 0／存在非 0」三段走齐.}
\bindp {\kaishu ③收束：证明题结论回代特位（通径、顶点、切位）复核；求参结果与前提（保持二次、Δ、范围）对账后再作答.}

% ============================================================
% 课堂评价 G5（恒5＝3单选+2填空＝导槽 G1~G5；ketang 新制顺排可跨栏，禁 \vbox 装栏）
% 印面号 38—42；多选门：本片正文多选 0 ≤4 合规
% ============================================================
\begin{ketang}{本组共5题（第38—42题），38—40为单选，41—42为填空．}

\jiancestem{{\fontsize{11.4pt}{14pt}\selectfont\heihao {\numboldjian 38}．}\hspace{4.1mm}\tieside{中档(知识点三)}\hspace{2.2mm}过抛物线 \(y^{2}=6x\) 的焦点作一条直线与抛物线交于 \(A(x_{1},y_{1})\)、\(B(x_{2},y_{2})\) 两点，若 \(x_{1}+x_{2}=3\)，则这样的直线的条数为\nobreak(\kongwei)}

\bindopt {\fontsize{10.5pt}{\qpTJgLeadPingjia}\selectfont \makebox[\dimexpr0.25\linewidth-0.5em\relax][l]{A．\(0\)}\makebox[\dimexpr0.25\linewidth-0.5em\relax][l]{B．\(1\)}\makebox[\dimexpr0.25\linewidth-0.5em\relax][l]{C．\(2\)}\makebox[\dimexpr0.25\linewidth-0.5em\relax][l]{D．\(3\)}\par}

{\ansblockgrayfalse
\begin{ansblock}[2章-导-课时17-G1]
% ans:2章-导-课时17-G1
\ansitem{38}{B}
\ansnote{详解}{焦点 \(F(3/2,0)\)．所求直线过焦点且与抛物线交于两点，可设 \(x=my+3/2\)（\(m\in R\)，\(m=0\) 即竖直线）．代入 \(y^{2}=6x\)：\(y^{2}-6my-9=0\)，\(\Delta=36m^{2}+36>0\) 恒成立，恒交于两点．\(x_{1}+x_{2}=m(y_{1}+y_{2})+3=m\cdot 6m+3=6m^{2}+3=3\)，得 \(m^{2}=0\)、\(m=0\)．故直线唯一：\(x=3/2\)，条数为 1，选 B．（验算：\(x=3/2\) 代入 \(y^{2}=6x\) 得 \(y^{2}=9\)，两交点 \(A(3/2,3)\)、\(B(3/2,-3)\)，\(x_{1}+x_{2}=3\) \ding{51}（此即通径位置）．）}
\end{ansblock}}

\jiancestem{{\fontsize{11.4pt}{14pt}\selectfont\heihao {\numboldjian 39}．}\hspace{4.1mm}\tieside{简单(知识点三)}\hspace{2.2mm}\(F_{1}\)、\(F_{2}\) 分别是双曲线 \(x^{2}/2-y^{2}/4=1\) 的左、右焦点，过 \(F_{1}\) 的直线分别交该双曲线的左、右两支于 \(A\)、\(B\) 两点，若 \(AF_{2}\perp BF_{2}\) 且 \(|AF_{2}|=|BF_{2}|\)，则 \(|AF_{2}|\) 等于\nobreak(\kongwei)}

\bindopt {\fontsize{10.5pt}{\qpTJgLeadPingjia}\selectfont \makebox[\dimexpr0.25\linewidth-0.5em\relax][l]{A．\(2\)}\makebox[\dimexpr0.25\linewidth-0.5em\relax][l]{B．\(2\sqrt{2}\)}\makebox[\dimexpr0.25\linewidth-0.5em\relax][l]{C．\(4\)}\makebox[\dimexpr0.25\linewidth-0.5em\relax][l]{D．\(4\sqrt{2}\)}\par}

{\ansblockgrayfalse
\begin{ansblock}[2章-导-课时17-G2]
% ans:2章-导-课时17-G2
\ansitem{39}{C}
\ansnote{详解}{\(a^{2}=2\)，\(2a=2\sqrt{2}\)．\(A\) 在左支：\(|AF_{2}|-|AF_{1}|=2a\)；\(B\) 在右支：\(|BF_{1}|-|BF_{2}|=2a\)．设 \(|AF_{2}|=|BF_{2}|=t\)，则 \(|AF_{1}|=t-2a\)、\(|BF_{1}|=t+2a\)．\(F_{1}(-\sqrt{6},0)\) 在左支顶点 \((-\sqrt{2},0)\) 的左方（位于左支凹侧内），过 \(F_{1}\) 的弦自 \(F_{1}\) 出发先遇左支点 \(A\)、再遇右支点 \(B\)，故 \(|AB|=|BF_{1}|-|AF_{1}|=4a=4\sqrt{2}\)．在 \(\mathrm{Rt}\triangle AF_{2}B\)（\(\angle AF_{2}B=90^{\circ}\)）中：\(|AF_{2}|^{2}+|BF_{2}|^{2}=|AB|^{2}\)，即 \(2t^{2}=(4\sqrt{2})^{2}=32\)，\(t^{2}=16\)，\(t=4\)（\(t>0\)）．故 \(|AF_{2}|=4\)，选 C．（验算：\(t=4\)：\(|AF_{1}|=4-2\sqrt{2}>0\) \ding{51}；\(|AB|=|BF_{1}|-|AF_{1}|=(4+2\sqrt{2})-(4-2\sqrt{2})=4\sqrt{2}\) \ding{51}；勾股 \(2\times 4^{2}=32=(4\sqrt{2})^{2}\) \ding{51}．）}
\end{ansblock}}

\jiancestem{{\fontsize{11.4pt}{14pt}\selectfont\heihao {\numboldjian 40}．}\hspace{4.1mm}\tieside{简单(知识点三)}\hspace{2.2mm}已知 \(F_{1}\)、\(F_{2}\) 是椭圆 \(x^{2}/16+y^{2}/9=1\) 的两焦点，过点 \(F_{2}\) 的直线交椭圆于 \(A\)、\(B\) 两点．在 \(\triangle ABF_{1}\) 中，若 \(|AB|+|AF_{1}|=10\)，则 \(|BF_{1}|\) 等于\nobreak(\kongwei)}

\bindopt {\fontsize{10.5pt}{\qpTJgLeadPingjia}\selectfont \makebox[\dimexpr0.25\linewidth-0.5em\relax][l]{A．\(2\)}\makebox[\dimexpr0.25\linewidth-0.5em\relax][l]{B．\(4\)}\makebox[\dimexpr0.25\linewidth-0.5em\relax][l]{C．\(6\)}\makebox[\dimexpr0.25\linewidth-0.5em\relax][l]{D．\(8\)}\par}

{\ansblockgrayfalse
\begin{ansblock}[2章-导-课时17-G3]
% ans:2章-导-课时17-G3
\ansitem{40}{C}
\ansnote{详解}{\(a=4\)，\(2a=8\)．\(F_{2}\) 在椭圆内部，弦 \(AB\) 过 \(F_{2}\)，故 \(F_{2}\) 在 \(A\)、\(B\) 之间，\(|AF_{2}|+|BF_{2}|=|AB|\)．\(\triangle ABF_{1}\) 的周长 \(|AB|+|AF_{1}|+|BF_{1}|=(|AF_{1}|+|AF_{2}|)+(|BF_{1}|+|BF_{2}|)=2a+2a=16\)（椭圆定义，与弦的位置无关）．由 \(|AB|+|AF_{1}|=10\) 得 \(|BF_{1}|=16-10=6\)，选 C．（验算：\(|BF_{1}|=6\) 落在焦半径范围 \([a-c,a+c]=[4-\sqrt{7},4+\sqrt{7}]\approx[1.35,6.65]\) 内 \ding{51}；周长恒等 \(4a\) 不依赖过焦弦的方向 \ding{51}．）}
\end{ansblock}}

\jiancestem{{\fontsize{11.4pt}{14pt}\selectfont\heihao {\numboldjian 41}．}\hspace{4.1mm}\tieside{简单(知识点三)}\hspace{2.2mm}过抛物线 \(T\)：\(y^{2}=2px\)（\(p>0\)）的焦点 \(F\) 的直线与 \(T\) 交于 \(A\)、\(B\) 两点，且 \(|AF|=2|FB|\)，\(T\) 的准线 \(l\) 与 \(x\) 轴交于点 \(C\)，若 \(\triangle CBF\) 的面积为 \(4\sqrt{2}\)，则 \(T\) 的通径长为\kongbai{}．}
\xiexwei{6mm}

{\ansblockgrayfalse
\begin{ansblock}[2章-导-课时17-G4]
% ans:2章-导-课时17-G4
\ansitem{41}{8}
\ansnote{详解}{设 \(A(x_{1},y_{1})\)、\(B(x_{2},y_{2})\)（\(y_{1}>0>y_{2}\)），直线 \(x=my+p/2\) 代入 \(y^{2}=2px\)：\(y^{2}-2pmy-p^{2}=0\)，得 \(y_{1}y_{2}=-p^{2}\)．由抛物线定义及 \(y_{1}^{2}=2px_{1}\)：\(|AF|=x_{1}+p/2=(y_{1}^{2}+p^{2})/(2p)\)，同理 \(|FB|=(y_{2}^{2}+p^{2})/(2p)\)．\(|AF|=2|FB|\) 即 \(y_{1}^{2}+p^{2}=2(y_{2}^{2}+p^{2})\)，得 \(y_{1}^{2}=2y_{2}^{2}+p^{2}\)．两边乘 \(y_{2}^{2}\) 并以 \((y_{1}y_{2})^{2}=p^{4}\) 代入：\(p^{4}=2y_{2}^{4}+p^{2}y_{2}^{2}\)．令 \(u=y_{2}^{2}\)：\(2u^{2}+p^{2}u-p^{4}=0\)，因式分解 \((2u-p^{2})(u+p^{2})=0\)，\(u>0\) 得 \(u=p^{2}/2\)，即 \(y_{2}^{2}=p^{2}/2\)（此时 \(y_{2}=-p/\sqrt{2}\)、\(y_{1}=\sqrt{2}p=-2y_{2}\) \ding{51}）．\(C(-p/2,0)\)、\(F(p/2,0)\)，\(|CF|=p\)，\(B\) 到 \(x\) 轴距离 \(|y_{2}|=p/\sqrt{2}\)，故 \(S_{\triangle CBF}=\frac{1}{2}\cdot p\cdot p/\sqrt{2}=p^{2}/(2\sqrt{2})=4\sqrt{2}\)，得 \(p^{2}=16\)、\(p=4\)，通径长 \(2p=8\)．（验算：\(p=4\)：\(C(-2,0)\)、\(F(2,0)\)、\(|y_{2}|=2\sqrt{2}\)，\(S=\frac{1}{2}\times 4\times 2\sqrt{2}=4\sqrt{2}\) \ding{51}；\(|AF|=(32+16)/8=6\)、\(|FB|=(8+16)/8=3\)，\(6=2\times 3\) \ding{51}．）}
\end{ansblock}}

\jiancestem{{\fontsize{11.4pt}{14pt}\selectfont\heihao {\numboldjian 42}．}\hspace{4.1mm}\tieside{简单(知识点一)}\hspace{2.2mm}过双曲线 \(C\)：\(x^{2}/2-y^{2}=1\) 上一点 \(P(2,1)\) 的切线方程是\kongbai{}．}
\xiexwei{6mm}

{\ansblockgrayfalse
\begin{ansblock}[2章-导-课时17-G5]
% ans:2章-导-课时17-G5
\ansitem{42}{y=x−1}
\ansnote{详解}{\(P(2,1)\)：\(4/2-1=1\) \ding{51}，\(P\) 在 \(C\) 上．设切线 \(l\)：\(y-1=k(x-2)\)，即 \(y=kx+(1-2k)\)．代入 \(x^{2}/2-y^{2}=1\)（乘 2）：\(x^{2}-2(kx+1-2k)^{2}=2\)，整理得 \((1-2k^{2})x^{2}-4k(1-2k)x-2[(1-2k)^{2}+1]=0\)．若 \(1-2k^{2}=0\)（\(k=\pm\sqrt{2}/2\)），方程降为一次，\(l\) 平行于渐近线 \(y=\pm(\sqrt{2}/2)x\) 而与 \(C\) 仅一个公共点，非切线，排除．故 \(1-2k^{2}\neq 0\)，\(l\) 与 \(C\) 相切即 \(\Delta=0\)：\par\noindent \(\Delta=16k^{2}(1-2k)^{2}+8(1-2k^{2})[(1-2k)^{2}+1]=0\)，除以 8 并展开（\((1-2k)^{2}=1-4k+4k^{2}\)）：\(2k^{2}(1-4k+4k^{2})+(1-2k^{2})(2-4k+4k^{2})=(2k^{2}-8k^{3}+8k^{4})+(2-4k+8k^{3}-8k^{4})=2k^{2}-4k+2=2(k-1)^{2}=0\)，\par\noindent 得 \(k=1\)（\(1-2\times 1=-2\neq 0\)，保持二次 \ding{51}）．故切线方程为 \(y=x-1\)．（验算：\(y=x-1\) 代入：\(x^{2}/2-(x-1)^{2}=1\) 即 \(-x^{2}/2+2x-2=0\)，亦 \(x^{2}-4x+4=0\)、\((x-2)^{2}=0\)，重根 \(x=2\)，切点即 \(P(2,1)\) \ding{51}——恰一公共点且不平行渐近线，为真切线；源解以导数求斜率系后置工具，本解改走联立 \(\Delta=0\)，读数一致．）}
\end{ansblock}}

\end{ketang}

% ---- 尾块（\tailfill 置 \end{multicols} 前；无参＝内容尾随框，件侧不擅自定高） ----
\tailfill

\end{multicols}
\end{document}
